% UTM.tex (RMCG20260406)

\input TM % \input RCstyle

\title{A Universal Turing Machine}
\author{Ramón Casares}
\contact{{\sc orcid:}
\URL 0000-0003-4973-3128<https://orcid.org/0000-0003-4973-3128>}
\subject{Computing}
\keywords{Turing completeness, universal Turing machine}

\maketitle

\beginabstract
We reconstruct the didactic universal Turing machine
by Marvin Minsky, MMo-UTM.
We test it on a Turing machine interpreter coded in Lua,
we comment on its limitations,
and we enhance it twice:
MMr-UTM,
 which can emulate any binary Turing machine; and
MMc-UTM,
 which can emulate any Turing machine.
The table of MMc-UTM is listed completely,
with comments and an example of emulation.
Since it has not any limitation
and it uses a straightforward encoding,
MMc-UTM makes a good constructive proof of
the theorem of Turing completeness.

\endabstract

\beginnote
This is \DOI{10.6084/m9.figshare.32261133}\new/.v4\wen/,
version {\tt\version}.\\
\copyright\ 2026 Ramón Casares; licensed as {\tt cc-by}.\\
Any comments on it to
{\tt\URL papa@ramoncasares.com<mailto:papa@ramoncasares.com>}
are welcome.
\endnote



\section{Introduction}

I defend that human language is a Turing complete language,
which is the language of a universal Turing machine.
Therefore, to me, the theorem of Turing completeness,
{\em there is a universal Turing machine}, is the key to language,
so I needed an accesible and foundational proof of it,
where foundational requires it to be constructive and self-contained.
And then, as we will see, I am afraid
I have had to be (a bit) picky on disclaimers.

Marvin \cite[\1 (\2  §7)] Minsky (1967) presented a UTM*,
that is, a universal Turing machine with a disclaimer.
We will call it MMo-UTM.
Disclaimer: MMo-UTM can only emulate
 Turing machines using two symbols, {\em blank} and another one,
 that work on a singly infinite tape, that is,
 MMo-UTM can only emulate binary and singly infinite Turing machines.
In particular, MMo-UTM cannot emulate itself.

We will discuss all MMo-UTM limitations
and we will proceed to eliminate them,
presenting two universal Turing machines in complete detail.
\beginpoints
\point Firstly, MMr-UTM, a revised version of MMo-UTM,
 which can emulate any binary Turing machine, but not itself.
\point And finally, MMc-UTM, which can emulate any Turing machine,
 so it can emulate itself.
 MMc-UTM is a universal Turing machine without disclaimers.
\endpoints

\break

In this paper, section \refsc{Computing},
we firstly introduce the Turing machine.
We start the section with an overview,
 in subsection \refsc{Turing machine},
and a more detailed definition, in subsection \refsc{Processor}.
Then we will see some examples, in subsection \refsc{Examples},
with the help of a Turing machine interpreter coded in Lua,
which we have developed for testing.
The section finishes, in subsection \refsc{Turing completeness},
 defining {\em Turing completeness} and {\em universal Turing machine}.

Next, in section \refsc{MMo-UTM}, we reconstruct MMo-UTM.
We start the section explaining its tape,
in subsection \refsc{MMo-Tape}, and,
in subsection \refsc{MMo-Table}, we recreate
 the \refx{diagram by Minsky that defines MMo-UTM}{MMoDiagram},
from which we obtain its table, which is listed in full.
Then, in subsection \refsc{MMo-Testing},
we test it using the Turing machine interpreter.
After doing some tests, we finish the section,
 in subsection \refsc{MMo-Limitations},
analyzing its limitations.

Then, in section \refsc{MMr-UTM},
we present its first enhancement, MMr-UTM,
a Turing machine that can emulate any binary Turing machine.
The first subsection, \refsc{MMr-Tape}, describes its tape,
compared with MMo-UTM's tape.
Then, in subsection \refsc{MMr-Table}, we present its
three new sub-machines, each one solving a limitation of MMo-UTM,
the patch that transforms MMo-UTM into MMr-UTM,
and the complete \refx{diagram of MMr-UTM}{MMrDiagram}.
In subsection \refsc{MMr-Testing}, we finish the section
testing MMr-UTM.

Section \refsc{Binary Turing completeness theorem}
is devoted to the theorem of binary Turing completeness,
which states that there is a Turing machine
that can emulate any binary Turing machine.
MMr-UTM is a constructive proof of the theorem.
This theorem together with another one by Shannon
imply that there is a universal Turing machine.

Therefore, in section \refsc{MMc-UTM},
we present a universal Turing machine.
Since it was designed on top of MMo-UTM
and it completely eliminates its limitations,
we call it MMc-UTM.
We start the section explaining its tape,
in subsection \refsc{MMc-Tape},
and then, in subsection \refsc{MMc-Table},
we present the whole table and
\refx{diagram of MMc-UTM}{MMcDiagram},
a comment about each of its sub-machines,
and a complete example of MMc-UTM emulating
a Turing machine. In subsection \refsc{MMc-Testing},
we finish the section testing MMc-UTM.

Section \refsc{Turing completeness theorem}
is devoted to the theorem of Turing completeness,
which says that there is a universal Turing machine,
which is a Turing machine that can emulate any Turing machine.
MMc-UTM is a constructive proof of the theorem.

In section \refsc{Discussion}, we comment
on the definition of universal Turing machine,
and on the merit of an independent and concrete
proof of the theorem of Turing completeness.
The first universal computing machine
was defined and designed by Turing,
and its limitations are reviewed in subsection \refsc{Turing}.
Then, in subsection \refsc{Encoding},
we present the problem of encoding,
raised by Davis and resolved by Minsky,
who requires to encode program and data separately.
In subsection \refsc{Definition},
we explain that Minsky's separation requirement
was the reason why
Turing's original definition of universal
Turing machine has been superseded.
Finally, in subsection \refsc{Foundations},
we argue that the proof of the theorem of
Turing completeness presented in this paper
is foundational,
because it is self-contained and constructive,
since it presents a concrete Turing machine,
MMc-UTM, which is universal without limitations
and uses a straightforward encoding.

We conclude, in section \refsc{Conclusion},
with a short summary.


\vfil\break

\section{Computing}

\subsection{Turing machine}

\dest
The seminal reference of computing is
its founding paper by \cite Turing (1936),
where he presents his $a$-machine,
 now known as Turing machine,
which is his mathematical model for computing devices.
A \define|Turing-machine|Turing machine|a tape controlled by a processor, $\cal P$|
 is a tape controlled by a processor.

\dest
The \define|tape||%
 inexhaustible chain of read and write squares|
is read and write memory, and it is divided in squares,
where each square can contain one symbol
out of a finite set of symbols, or none.
When a square contains no symbol, we say that it is
\define|blank||empty tape square
 and its corresponding special symbol|.
We will consider that {\em blank} is a symbol,
 though a {\em special} one.
Each square has exactly one neighbor square to the left
and one neighbor square to the right, meaning
that the tape is infinite and
that no computation will abort because of a lack of squares,
though the number of non-{\em blank} squares will
always be finite. Another way of seeing it is
that the tape is finite but potentially infinite, and
that new {\em blank} squares are created
whenever either end of the tape is going to be overstepped.

\dest
The \define|processor||a finite-state automaton|
is a finite-state automaton.
Then, the processor is in
each moment in exactly one state out of a finite set of states,
and the next state and the output depend on the current
state and the input.
In the case of the Turing machine processor,
the input is read from the scanned square of the tape, and
the output is a pair composed of a symbol and a movement,
where the symbol is written to the scanned square and
the \define|movement||next scanned square,
 either {\em left} or {\em right}|
is {\em left} or {\em right},
so the next scanned square will be
the neighbor one to the left,
or the neighbor one to the right.
So a {\em table} defining the next state
and the pair of symbol to write and movement to do
for each possible current state and read symbol
determines completely the Turing machine behavior.
We will say that the pair of current state ({\sc c})
and read symbol ({\sc r}) is the
\define|index||instruction reference:
 current state ({\sc c}) and read symbol ({\sc r})|,
and that the triplet of next state ({\sc n}),
symbol to write ({\sc w}) and movement to do ({\sc m})
is its corresponding
\define|instruction||what to do: either next state ({\sc n}),
 write symbol ({\sc w}) and movement ({\sc m}) or {\em halt}|.
Together, the index with its corresponding instruction
form a \define|quintuple||index and corresponding instruction
 specifying a transition|,
which specifies a \define|transition||a computing step
 specification or what to do in a specific situation|.
As the number of states and the number of symbols are both finite,
the table is finite.

\dest
The processor state, the tape content and the scanned square
define the \define|situation||current state,
 tape content and scanned square| of the Turing machine.
A \define|computation||from the initial situation,
 follow instructions till {\em halt}, if ever|
starts from a well-defined situation and is executed in a loop
with four stages ---read, transition, write, and move---
that is exited whenever the transition is a halting one,
and then we say that the Turing machine has halted.
We will assume without loss of generality that
in the initial situation the state is the
\define|initial state||state when starting a computation|
and the scanned square is any one such that
all squares to its left are {\em blank},
preferably the leftmost non-{\em blank} square.
From that initial situation, the Turing machine
follows the instructions given in the table
until the instruction to do is {\em halt},
thus ending the computation, or else it goes forever.
Some more details and examples are provided below,
in subsections \refsc{Processor} and \refsc{Examples}.

\dest
So the Turing machine transforms the string of symbols that was
written on the tape when it started into the string of symbols
that was written on the tape when it halted.
From this point of view, each Turing machine
implements a function from the set of strings of symbols
to the set of strings of symbols,
where the strings are finite in length and
its symbols are drawn from a finite set.
By definition, any function implemented by a Turing machine is
a \define|computable function||one implemented by a Turing machine|.
Note that not all these functions are total functions,
because there is not any guarantee that
a Turing machine on an initial string will halt,
and then the function implemented by the Turing machine
is not defined for those input strings on which the
machine never halts.


\subsection{Processor}

\dest
A Turing machine $\cal P$ processor has four main components:\\
the states $\Phi$, the symbols $\Gamma$,
the movements $\Delta$, and the table $\Pi$.
\beginpoints
\point The set of \define|states||processor’s set of states, $\Phi$|
 $\Phi$ has to be finite, and
 the {\em initial} state, let us call it $\M A$, has to be in it,
 so $\M A \in\Phi$ and $|\Phi| \ge 1$.
\point The set of \define|ordinary symbols||non-empty set of
 non-{\em blank} symbols, $\bar\Gamma$|
 $\bar\Gamma$ has to be finite and not empty. The set of
 \define|symbols||tape’s set of symbols,
 $\Gamma = \bar\Gamma \cup \{\,\hbox{\it blank}\,\}$, or alphabet|
 $\Gamma$ is set $\bar\Gamma$ plus the {\em blank} special symbol,
 let us write it \M0, that denotes that a tape square is {\em blank}.
 Then, $\Gamma = \bar\Gamma \cup \{ \M0 \}$,
 so $\M0\in\Gamma$ and $|\Gamma| \ge 2$.
 We say that $\Gamma$ is the
 \define|alphabet||well-ordered set of symbols, $\Gamma$|
 to say that its symbols are well-ordered.
\point The set of \define|movements||set of movements,
 $ \Delta=\{\, \hbox{\it left}, \hbox{\it right}\, \}$|
 $\Delta$ has two elements:
 {\em left}, let us write it \M<; and
 {\em right}, let us write it \M>.
 Then, $\Delta=\{\M<, \M>\}$.
\point The \define|table||partial function
 $\Pi: \Phi\times\Gamma \to \Phi\times\Gamma\times\Delta$|
 $\Pi$ is a partial function
 $\Pi: \Phi\times\Gamma \to \Phi\times\Gamma\times\Delta$,
 so it can be written as a table
 of $|\Phi|\times|\Gamma|$ rows at most, and five columns.
 Function $\Pi$ argument is the {\em index}, and
 its value is its corresponding {\em instruction}.
\endpoints

\dest
The default instruction is
\define|halt||instruction that finishes a computation|.
The {\em halt} instruction finishes the computation, and
it does not change anything, neither the state nor the symbol,
and it does not move.
Being the default instruction means that,
 for any argument of table $\Pi$ that is not defined,
 the corresponding instruction is {\em halt},
 so every argument gets its corresponding instruction.
 %%%Remember: any row omitted from table~$\Pi$
 %%%specifies a halting transition.
This way, table $\Pi$ with the default instruction {\em halt} makes
the \define|transition function||total function
 $ \Pi' : \Phi\times\Gamma \to
 (\Phi\times\Gamma\times\Delta)\cup\{\,\hbox{\it halt}\,\}$|
 $\Pi'$ a total function,
$\Pi':  \Phi\times\Gamma \to
(\Phi\times\Gamma\times\Delta)\cup\{\,\hbox{\it halt}\,\}$.

Set $\Gamma^*$ is the infinite set
 of the {\em finite} strings drawn from $\Gamma$.
There is a practical problem regarding the {\em blank} special symbol,
because there are two infinities of them on the tape,
one to the left and the other to the right.
We will ignore these two infinities,
but we should be aware that
a meaningful finite string of {\em blank} symbols
concatenated with any of those two infinities could
cause ambiguities.

\dest
We use \define|meta-string infty|meta-string $\infty$|a
 computation that does not halt|
 $\notin \Gamma^*$ to denote that a computation does not halt.
Overloading the Turing machine name, for instance $\cal P$,
which can refer to the Turing machine, to its processor,
or to the implemented function,
we will write $\TM P<d> = \TM<r>$ to express that,
when we write string ${\frak d}\in\Gamma^*$
on the tape of Turing machine $\cal P$,
it computes it and halts
 with the string ${\frak r}\in\Gamma^*$ on the tape.
So $\frak d$ is the input string to $\cal P$,
and $\frak r$ is the corresponding output string.
Should it no halt, we would write $\TM P<d> = \infty$ instead,
expressing that function $\cal P$ is not defined
for argument $\frak d$.

\dest\new/\wen/
We use \define|meta-symbol *|meta-symbol {\tt *}|any other symbol in index,
 and the same symbol in instruction|
 $\notin \Gamma$ to denote {\em any other symbol}.
That is, index \M{C*} applies when, in state \M{C},
the read symbol is any not explicitly defined.
And, in that case, the corresponding instruction
can also use a \M*, then meaning {\em the same symbol},
which is the read one.
For example, if index \M{C1} is not in the table,
then row \M{C*B*<} expands to \M{C1B1<}.
Meta-symbol \M* is handy when writing tables,
 but it can always be replaced by symbols;
meta-symbol \M* is just an abbreviation.

\dest
There are many versions of the Turing machine,
and this is only one of them,
basically the one by \cite Minsky (1967).
Some use quadruples instead of quintuples,
so the instruction is a pair composed of
the next state and the action,
where the action can either be to write a symbol or to move.
Others halt when reaching a halting state;
we will use them frequently, since for us
a \define|halting state||one from which all
 transitions are halting ones|
 is any from which any transition is a halting transition.
In other words, a halting state
can be completely omitted from the {\em index}.
Some others require two halting states,
the accepting and the rejecting state,
which we can also implement,
and even more, since we are not limited to two halting states.


\vfill\break

\subsection{Examples}

For example, this is the table for a Turing machine ${\cal S}_2$
that implements the successor function in base two without zero.
In order to avoid ambiguities,
we will not use the {\em blank} symbol to write the numbers,
and then we will need a three-symbol alphabet; we will use
alphabet $\Gamma=\{ \M0, \M1, \M2 \}$, so $|\Gamma| = 3$.
We have added a header to help in reading each quintuple:
{\sc c} for the current state,
{\sc r} for the read symbol,
{\sc n} for the next state,
{\sc w} for the written symbol, and
{\sc m} for the movement.
Four states and three symbols make a twelve-row table,
but three are omitted because there are three halting transitions;
in fact, state \M{H} is a halting state
because all transtions from state \M{H} are halting ones.
$${\cal S}_2\quad\Table{
 A& 0& A& 0& >\cr
 A& 1& B& 1& >\cr
 A& 2& B& 2& >\cr
 B& 0& C& 0& <\cr
 B& 1& B& 1& >\cr
 B& 2& B& 2& >\cr
 C& 0& H& 1& <\cr
 C& 1& H& 2& <\cr
 C& 2& C& 1& <\cr
}$$

To see how this specific Turing machine ${\cal S}_2$ works,
we can see what happens if we write the string \M{2212} on the tape,
where the numeral \M{2212} in base two without zero
is number twenty eight ($\M2.2^3+\M2.2^2+\M1.2^1+\M2.2^0$).
That is, we will write the four symbols \M2, \M2, \M1, and \M2
from left to right in four consecutive squares.
Then, the initial situation will be as follows.
\T{0~0~0~A2~2~1~2~0~0~0}\\
There are two infinities of {\em blank}s,
one to the left and one to the right, which can be omitted,
though we have kept three {\em blank}s
on each side for didactic purposes.
The scanned square is the one that has attached to it
the current state, which is \M{A} at the start.
Now, \M{A2} is the index of the third row in the table,
so the instruction to do is \M{B2>},
and therefore the next situation,
after going to state \M{B}, rewriting the \M2, and moving to the right,
is as follows.
\T{0~0~0~2~B2~1~2~0~0~0}\qquad
For index \M{B2}, the instruction is also \M{B2>}.
\T{0~0~0~2~2~B1~2~0~0~0}\qquad
Now the index is \M{B1}, so the instruction is \M{B1>}.
\T{0~0~0~2~2~1~B2~0~0~0}\qquad
And, again, the instruction for index \M{B2} is \M{B2>}.
\T{0~0~0~2~2~1~2~B0~0~0}\qquad
For \M{B0}, it is \M{C0<},
so the machine changes the state.
\T{0~0~0~2~2~1~C2~0~0~0}\qquad
And, for index \M{C2}, the instruction is \M{C1<}.
\T{0~0~0~2~2~C1~1~0~0~0}\qquad
Now, the instruction for \M{C1} is \M{H2<}.
\T{0~0~0~2~H2~2~1~0~0~0}\qquad\dest
Since index \M{H2} is not in the table,
the instruction to do is {\em halt}, so
we have reached the final situation.
Then the result of the computation,
from which the two infinities of {\em blank} symbols are omitted,
is the string \M{2221}, which is twenty nine in base two without zero.
We summarize it, thus:
$\RS S_{\rm 2}<2212> = \RS<2221>$.
Defining \define|step||transition from a computing situation
 to the next one|
 as a transition from a computing situation to the next one,
this computation takes seven steps.
Now you should try yourself
$\RS S_{\rm 2}<22> = \RS<111>$,
and others.

\vfil\break

And, after having tried yourself some computations,
you can verify whether you are doing it rigthly or wrongly
using the Turing machine interpreter in GitHub.
Clone the repository
{\tt\URL github.com/ramoncasares/UTM<https://github.com/ramoncasares/UTM.git>},
go to its dir "lua",
and execute the Lua interpreter in inter\-active mode,
which uses \M> as prompt.
Firstly, we will ask the interpreter to redo
the previous computations.
\begincode
...$ git clone https://github.com/ramoncasares/UTM.git
...$ cd UTM/lua/
.../UTM/lua$ lua -i
Lua 5.4.7  Copyright (C) 1994-2024 Lua.org, PUC-Rio
> mac = require("machine")
> TMS2t = "A0A0> A1B1> A2B2> B0C0< B1B1> B2B2> C0H1< C1H2< C2C1<"
> TMS2 = mac:new(TMS2t)
> TMS2:tape("2212")
2212
> TMS2:compute(true)
2221   |A2212|->|2B212|->|22B12|->|221B2|->|2212B0|->|221C20|->|22C110|->|2H2210|
> TMS2:tape("22")
22
> TMS2:compute(true)
111   |A22|->|2B2|->|22B0|->|2C20|->|C210|->|C0110|->|H01110|
\endcode
And we can continue.
Calling "compute" without arguments, it returns the tape
and, instead of the steps, the number of steps \new/executed\wen/.
\new/Function "compute" resets the Turing machine to the initial situation,
so it starts from the initial state and
scanning the leftmost non-{\em blank} square, if any,
but it does not change the tape.\wen/
\begincode
> TMS2:compute()
112   5
> TMS2:compute()
121   6
> TMS2:compute()
122   5
\endcode
\new/However, since Lua function "compute" will not return (gracefully)
if the corresponding computation does not halt,
we can rather call function "steps", which always returns.
\label{steps}Function "tm:steps(n,...)"\  runs Turing machine "tm"
 from its current situation, that is,
 neither resetting it nor changing its tape,
 until some condition is met,
 and returns the new situation and the number of steps executed.
If Turing machine "tm" halts, then the function returns.
It also returns after having executed "n" steps.
Each of the other optional arguments ("...")\ refers
either to a state "s" (a "string"),
 and then the function returns if the current state is "s",
or to a transition "{st,sy}" (a state and symbol "table"),
 and then it returns
 if the current state is "st" and the scanned symbol is "sy".\wen/
For more details on these and other functions,
you can examine the source code,
which is in directory ".../UTM/lua".
\new/\begincode
> TMS2:tape("")
0
> TMS2:steps(10)
0000000000[A]0   10
> TMS2:steps(10)
00000000000000000000[A]0   10
> TMS2:steps(10)
000000000000000000000000000000[A]0   10
> os.exit()
\endcode
\wen/

\vfil\break

We could have saved some ink using the meta-symbol \M*.
Both represent the very same table of Turing machine ${\cal S}_2$,
although the following one requires less typing.
$${\cal S}_2\quad\Table{
 A& 0& A& 0& >\cr
 A& *& B& *& >\cr
 B& 0& C& 0& <\cr
 B& *& B& *& >\cr
 C& 0& H& 1& <\cr
 C& 1& H& 2& <\cr
 C& 2& C& 1& <\cr
}$$

It is perhaps a bit easier for us to read this reduced form.
It makes clearer that, when ${\cal S}_2$ is in state \M{A},
it will keep moving right while reading {\em blank} symbols,
and that it will change to state \M{B}
when it finds a non-{\em blank} one, that is, an ordinary symbol.
This is the standard way to implement the convention that,
 in the initial situation, the scanned square is
the leftmost one where an ordinary symbol is written,
 or any to its left. \dest
This convention makes convenient to interpret the
state of the index of the first row of the table
 to be the {\em initial} state,
the symbol of the index of the first row of the table
 to be the {\em blank} symbol, and
the movement of the instruction of the first row of the table
 to be the {\em right} movement.
Whenever possible, we will take advantage of this
\define|first row rule||the first row declares the {\em initial} state,
 the {\em blank} symbol and the {\em right} movement|,
by which the table contains all the information
that defines a Turing machine.

We can use the interpreter to verify that it is the same Turing machine.
\begincode
.../UTM/lua$ lua -i
Lua 5.4.7  Copyright (C) 1994-2024 Lua.org, PUC-Rio
> mac = require("machine")
> TMS2xt = "A0A0>A*B*>B0C0<B*B*>C0H1<C1H2<C2C1<"
> TMS2x = mac:new(TMS2xt)
> TMS2x:tape("2212")
2212
> TMS2x:compute(true)
2221   |A2212|->|2B212|->|22B12|->|221B2|->|2212B0|->|221C20|->|22C110|->|2H2210|
> TMS2x:tape("22")
22
> TMS2x:compute(true)
111    |A22|->|2B2|->|22B0|->|2C20|->|C210|->|C0110|->|H01110|
> TMS2x:compute()
112   5
> Ctrl-D
\endcode

\vfil\break

Another simple Turing machine,
which we will use later, is "TM1022".
$$\M{TM1022}\quad\Table{
 A& 0& B& 1& >\cr
 A& 1& A& 1& <\cr
 B& 0& C& 2& >\cr
 B& 1& H& 0& <\cr
 B& 2& A& 2& <\cr
 C& 0& B& 2& <\cr
}$$
This table of Turing machine "TM1022"
honors the first row rule,
and then
the {\em initial} state is "A",
{\em blank} is "0",
and {\em right} is ">".

When we run \M{TM1022} on a blank tape,
meaning that all its squares are {\em blank}, we get \M{1022},
where we have omitted the infinities of {\em blank}s.
We can verify it using the interpreter.
\begincode
.../UTM/lua$ lua -i
Lua 5.4.7  Copyright (C) 1994-2024 Lua.org, PUC-Rio
> mac = require("machine")
> TM1022 = mac:new"A0B1> A1A1< B0C2> B1H0< B2A2< C0B2<"
> TM1022:compute(true)
1022    |A0|->|1B0|->|12C0|->|1B22|->|A122|->|A0122|->|1B122|->|H1022|
> os.exit()
\endcode
All these examples and other ones are in file "test/machine.lua",
which is in directory "test".
To see them, you should go to "test", set "LUA_PATH",
and read from "machine.lua".
\begincode
.../UTM/lua$ cd ../test/
.../UTM/test$ export LUA_PATH="../lua/?.lua"
.../UTM/test$ cat machine.lua - | lua -i
Lua 5.4.7  Copyright (C) 1994-2024 Lua.org, PUC-Rio
> -- test/machine.lua (RMCG20260116)
>
> mac = require("machine")
>
> -- TM1022 on a blank tape returns 1022
> TM1022 = mac:new"A0B1> A1A1< B0C2> B1H0< B2A2< C0B2<"
> TM1022:compute(true)
1022    |A0|->|1B0|->|12C0|->|1B22|->|A122|->|A0122|->|1B122|->|H1022|
> TM1022:compute(true)
10022   |A1022|->|A01022|->|1B1022|->|H10022|
> TM1022:compute(true)
100022  |A10022|->|A010022|->|1B10022|->|H100022|
.
.
.
> -- end test/machine.lua
> TM1022:compute(true)
1000022  |A100022|->|A0100022|->|1B100022|->|H1000022|
> TM1022:compute(true)
10000022 |A1000022|->|A01000022|->|1B1000022|->|H10000022|
\endcode

\vfil\break

\label{TM1022b}
The set of states and the alphabet of Turing machine "TM1022" are:
$\Phi = \{ \M{A}, \M{B}, \M{C}, \M{H} \}$ and
$\Gamma = \{ \M{0}, \M{1}, \M{2} \}$,
meaning that, in fixed length binary,
two bits are needed to encode the states and
two bits are needed to encode the symbols.
Applying the finite map
$$\{ \M{A}\rightarrow\M{00}, \M{B}\rightarrow\M{01},
     \M{C}\rightarrow\M{10}, \M{H}\rightarrow\M{11};\;
     \M{0}\rightarrow\M{00},
     \M{1}\rightarrow\M{10}, \M{2}\rightarrow\M{11};\;
     \M{<}\rightarrow\M{0},  \M{>}\rightarrow\M{1}\}$$
we get the binary encoded Turing machine "TM1022b".
$$\M{TM1022b}\quad\Table{
 00& 00& 01& 10& 1\cr
 00& 10& 00& 10& 0\cr
 01& 00& 10& 11& 1\cr
 01& 10& 11& 00& 0\cr
 01& 11& 00& 11& 0\cr
 10& 00& 01& 11& 0\cr
}$$
Running \M{TM1022b} on a blank tape we get \M{10 00 11 11},
where we have omitted the infinities of {\em blank}s and
each symbol is coded by two bits.
Applying the map for the alphabet in reverse
$$\{  \M{0}\leftarrow\M{00},
     \M{1}\leftarrow\M{10}, \M{2}\leftarrow\M{11} \}$$
we get "1022", which is what "TM1022" computed.
We can verify it using the interpreter.
\begincode
> code = require("code")
> TM1022b = code.tobinary(TM1022)
> TM1022b:tape("")
00
> TM1022b:compute({prest = "[", postst = "]"})
10 00 11 11        |[00] 00|->|10 [01] 00|->|10 11 [10] 00|->|10 [01] 11 11|->\
   |[00] 10 11 11|->|[00] 00 10 11 11|->|10 [01] 10 11 11|->|[11] 10 00 11 11|
> Ctrl-D
\endcode


\subsection{Turing completeness}

\leavevmode
\define|Turing-completeness|Turing completeness|computing
 can express computing completely|
is a property of computing, namely,
that {\em computing can express computing completely}.
The original formulation is in \cite Turing (1936), page 241,
in the first sentence of section §6,
titled {\sl The universal computing machine}:
“It is possible to invent a single machine
which can be used to compute any computable sequence.”

In modern terms, a \define|universal Turing machine||a
 Turing machine $\cal U$ that can emulate any Turing machine|
$\cal U$ is a Turing machine that can emulate
 any Turing machine $\cal P$.
The emulation has to be such that,
if a computation by $\cal P$ halts and computes a string,
then the emulation by $\cal U$ also halts and computes the same string,
and if a computation by $\cal P$ does not halt,
then the emulation by $\cal U$ does not halt either.

Near the end of this paper, in \refsc{Discussion},
we will see why the modern formulation superseded
the original one by Turing. But, before that,
 in order to understand Turing completeness better,
we will see three universal Turing machines in complete detail.
So, next, we will present and reconstruct
the didactic universal Turing machine
designed by \cite Minsky (1967), §7.
Here we will call it MMo-UTM.


\vfill\break

\section{MMo-UTM}

\labeled{MMo-Tape}\subsection{Tape}

In order to emulate a binary Turing machine ${\cal P}_2$,
MMo-UTM uses the eight-symbol alphabet
$\Gamma_{\!\rm o} = \{ \M0, \M1, \M A, \M B, \M M, \M S, \M X, \M Y \}$,
where "0" is the {\em blank} symbol.
The initial tape only contains symbol "M" once,
symbol "Y" twice, symbol "X" as many times as quintuples in ${\cal P}_2$,
symbol "1" a finite number of times, and any number of {\em blank}s, "0".
MMo-UTM's tape has three areas:
\beginpoints
\point {\em tape}: leftly infinite emulated tape of ${\cal P}_2$,
 where the symbol "M" marks the emulated scanned square;
\point {\em index}: current/next state and
 read/write symbol of ${\cal P}_2$;
\point {\em program}: description of ${\cal P}_2$,
 that is, its table (free of meta-symbols) coded in fixed length binary,
 so each state is a different binary string of length "s",
 and linearized, which is done by prefixing each quintuple with an "X".
 The {\em blank} of ${\cal P}_2$ is coded \M0,
 its other symbol is coded \M1;
 {\em left} is coded \M0, and {\em right} is coded \M1.
\endpoints

A valid initial tape is as follows,
where value of the bit in the "M" square
 is the value of rightmost bit in the {\em index},
and the {\em initial} state is the one written
 in the  {\em state} area of {\em index}.
$$\catcode`\<=1 \catcode`\>=2 \catcode`\{=12 \catcode`\}=12 \tt
\hbox<0*[01]*M[01]*Y[01]{s}[01](X[01]{s}[01][01]{s}[01][01])+Y0*>
$$
Our regular expressions use the following meta-symbols:
\beginpoints
\point postfix "*", which means zero or more times;
\point postfix "+", which means one or more times;
\point postfix "?", which means optional, that is, either once or none;
\point postfix "{n}", which means "n" times;
\point the square brackets "[" and "]", which mean
 any of the symbols between them,\\
 so "[01]" means either "0" or "1"
 and "[01]{s}" stands for any binary string of length "s",
 where "s" is the number of bits that encode
 each state of ${\cal P}_2$; and
\point  the parentheses "(" and ")", which enclose regular expressions.
\endpoints

\MT:
\MTbeginfig(15cm,3cm,0pt);
\MT: pickup medpen;
%\MT: draw (0,0) -- (w,0) -- (w,h) -- (0,h) -- cycle;
\MT: numeric cn; cn = 64; numeric ch; ch = 20pt;
\MT: draw (0,h) .. (w,h);
\MT: draw (0,h-ch) .. (w,h-ch);
\MT:
\MT: for n := 1 step 2 until cn:
\MT:  draw (n*w/cn,h) .. (n*w/cn,h-ch); endfor
\MT: string tape;
\MT: tape = "00M11000Y010X0101110X0110110Y00";
\MT: string sy;
\MT: defaultfont:="cmtt12";
\MT: for n := 0 step 1 until (length tape):
\MT:  sy := substring (n,n+1) of tape;
\MT:  z[n].lbl = ( 2*(n+1)*w/cn, h-ch/2 );
\MT:  label(sy,z[n].lbl);
\MT: endfor
\MT: pickup thinpen;
\MT: defaultfont:="cmti12";
\MT: labeloffset:=3pt;

\MT: draw (17w/cn,h-ch) .. (17w/cn,0);
\MT: drawarrow  (17w/cn,8pt) .. (0,8pt);
\MT: label.top("tape",(8.5w/cn,8pt));

\MT: draw (19w/cn,h-ch) .. (19w/cn,0);
\MT: draw (25w/cn,h-ch) .. (25w/cn,0);
\MT: drawdblarrow  (19w/cn,8pt) .. (25w/cn,8pt);
\MT: label.top("index",(22w/cn,8pt));

\MT: draw (23w/cn,h-ch) .. (23w/cn,(h-ch)/2);
\MT: drawdblarrow  (19w/cn,(h-ch)/2+8pt) .. (23w/cn,(h-ch)/2+8pt);
\MT: label.top("state",(21w/cn,(h-ch)/2+8pt));

\MT: draw (27w/cn,h-ch) .. (27w/cn,(h-ch)/2);
\MT: draw (41w/cn,h-ch) .. (41w/cn,(h-ch)/2);
\MT: drawdblarrow  (27w/cn,(h-ch)/2+8pt) .. (41w/cn,(h-ch)/2+8pt);
\MT: label.top("quintuple",(34w/cn,(h-ch)/2+8pt));

\MT: draw (43w/cn,h-ch) .. (43w/cn,(h-ch)/2);
\MT: draw (57w/cn,h-ch) .. (57w/cn,0pt);
\MT: drawdblarrow  (43w/cn,(h-ch)/2+8pt) .. (57w/cn,(h-ch)/2+8pt);
\MT: label.top("quintuple",(50w/cn,(h-ch)/2+8pt));

\MT: drawdblarrow  (25w/cn,8pt) .. (57w/cn,8pt);
\MT: label.top("program",(41w/cn,8pt));

\MTendfig;

\bigskip
\centerline{\box\MTbox}
\bigskip


\labeled{MMo-Table}\subsection{Table}

In the next page, you can find a recreation of
Fig.~7.2-9, page 142, in \cite Minsky (1967).
We have named the states \M{qmn},
where \M{m} is a sub-machine number, and \M{n} a state number.
We use \M{<} for the left movement, instead of Minsky's \M{L},
and \M{>} for the right movement, instead of his \M{R}.

\vfill\break

\MT:
\MTbeginfig(15.5cm,23.5cm,0pt); %% MMo-UTM Diagram
\MTresetpa;
\MT: pickup medpen;
\MT: defaultfont:="cmtt12";
\MT: draw (0,0)--(w,0)--(w,h)--(0,h)--cycle withcolor green;
\MTcode
 z99 = (6w/8, 9h/11);
 z11 = (6w/8, 6.5h/11);
 z12 = (2w/8, 8h/11);
 z13 = (3w/8, 9.5h/11);
 z14 = (3w/8, 6.5h/11);
 z15 = (4w/8, 8h/11);
 z16 = (6w/8, 8h/11);
 z21 = (2w/8, 6h/11);
 z22 = (3w/8, 5h/11);
 z23 = (1w/8, 5h/11);
 z24 = (3w/8, 4h/11);
 z25 = (1w/8, 4h/11);
 z26 = (2w/8, 3h/11);
 z31 = (3w/8, 2h/11);
 z32 = (1w/8, 2h/11);
 z33 = (2w/8, 1h/11);
 z34 = (4w/8, 1h/11);
 z35 = (6w/8, 1h/11);
 z41 = (5w/8, 3h/11);
 z42 = (7w/8, 3h/11);
 z43 = (5w/8, 4h/11);
 z44 = (7w/8, 4h/11);
 z45 = (5w/8, 5h/11);
 z46 = (7w/8, 5h/11);

\MTline{ label("HALT",(z99+(0,24pt))) withcolor red;}
\MTstate(99)"q99"">";
\MTline{ label("START",(z11+(-24pt,24pt))) withcolor blue;}
\MTstate(11)"q11""<";
\MTstate(12)"q12"">";
\MTstate(13)"q13"">";
\MTstate(14)"q14"">";
\MTstate(15)"q15""<";
\MTstate(16)"q16"">";

\MTstate(21)"q21"">";
\MTstate(22)"q22""<";
\MTstate(23)"q23""<";
\MTstate(24)"q24"">";
\MTstate(25)"q25"">";
\MTstate(26)"q26"">";

\MTstate(31)"q31""<";
\MTstate(32)"q32""<";
\MTstate(33)"q33"">";
\MTstate(34)"q34"">";
\MTstate(35)"q35""<";

\MTstate(41)"q41""<";
\MTstate(42)"q42""<";
\MTstate(43)"q43""<";
\MTstate(44)"q44"">";
\MTstate(45)"q45"">";
\MTstate(46)"q46"">";

\MT: pickup thinpen;

\MTauto(11)(viii,x)"0""A";
\MTauto(11)(iv,ii)"1""B";
\MTline{ z11.Y.p = (7w/8,9h/11);}
\MTline{ z11.Y.q = (6w/8,10h/11);}
\MTline{ z11.Y.t = (1w/8,9h/11);}
\MTinput(11.i)(12.x)"Y""Y";
\MTinput(12.i)(13.vii)"A""0";
\MTinput(12.v)(14.xi)"B""1";
\MTinput(12.vi)(21.xii)"X""X";
\MTinput(13.v)(15.xi)"0""A";
\MTinput(13.iv)(16.x)"1""B";
\MTinput(14.i)(15.vii)"1""B";
\MTinput(14.ii)(16.viii)"0""A";
\MTinput(15.ix)(12.iii)"Y""Y";
\MTinput(16.vi)(11.xii)"X""X";
\MTinput(16.xii)(99.vi)"Y""Y";

\MTinput(21.v)(22.xi)"0""A";
\MTinput(21.vii)(23.i)"1""B";
\MTinput(22.vi)(24.xii)"Y""Y";
\MTinput(23.vi)(25.xii)"Y""Y";
\MTinput(24.viii)(26.i)"0""A";
\MTinput(24.vii)(26.ii)"1""A";
\MTinput(24.vi)(31.xii)"X""X";
\MTinput(25.iv)(26.xi)"1""B";
\MTinput(25.v)(26.x)"0""B";
\MTinput(25.vi)(32.xii)"X""X";
\MTinput(26.xii)(21.vi)"X""X";

\MTinput(31.vii)(33.i)"M""A";
\MTinput(32.v)(33.xi)"M""B";
\MTauto(33)(ix,vii)"B""1";
\MTauto(33)(v,vii)"A""0";
\MTinput(33.iii)(34.ix)"X""X";
\MTinput(34.ii)(35.x)"0""0";
\MTinput(34.iv)(35.viii)"1""1";
\MTauto(35)(iii,v)"B""1";
\MTauto(35)(vii,v)"A""0";
\MTinput(35.xi)(41.v)"0""S";
\MTinput(35.i)(42.vii)"1""S";

\MTinput(41.xii)(43.vi)"A""0";
\MTinput(41.ii)(44.viii)"B""0";
\MTinput(42.x)(43.iv)"A""1";
\MTinput(42.xii)(44.vi)"B""1";
\MTinput(43.xii)(45.vi)"0""M";
\MTinput(43.ii)(46.viii)"1""M";
\MTinput(44.x)(45.iv)"0""M";
\MTinput(44.xii)(46.vi)"1""M";
\MTinput(45.i)(11.vii)"S""A";
\MTinput(46.xi)(11.v)"S""B";

\MTline{ label("MMo-UTM",(4w/8,12pt));}

\MTendfig;

\dest
\centerline{\define*|MMoDiagram|MMo-UTM|didactic universal
 Turing machine by Minsky, see \refsc{MMo-UTM}|\box\MTbox}
\vfill\break

This is the comma-separated values table of MMo-UTM,
as read from the diagram.
\doublecolumns
\file{../csv/MMo-UTM.csv}
\singlecolumn

MMo-UTM is a 8-symbol 24-state Turing machine.
One of the 24 states is the halting state \M{q99}.

This table does not obey the first row rule completely,
so we have to specify that the {\em right} movement is \M>,
and not the movement of the first row.
Since, in the {\em index} of the first row,
the state is the {\em initial} state \M{q11},
and the symbol is the {\em blank} symbol \M0,
these two do not need to be specified.

Sub-machines Locate and Copy implement the {\em transition} stage
of the computing loop, seen in \refsc{computation}.
The other two sub-machines implement
the other stages of the computing loop, but somehow mixed.
Since the initial state is \M{q11}, then MMo-UTM starts
when the {\em read} stage was already done.

\vfil\break


\labeled{MMo-Testing}\subsection{Testing}

Now, let us see MMo-UTM emulating "TM101",
which is a simple Turing machine that,
when it starts on a blank tape,
it halts with the string "101" on the tape.
$$\hbox{\tt TM101}\quad\Table{
 A& 0& B& 1& >\cr
 B& 0& C& 0& >\cr
 C& 0& H& 1& <\cr
}$$
These are the first commands of "test/MMo.lua", a Lua session.
Do not forget setting \verb|LUA_PATH|
when working from a different directory.
\begincode
.../UTM/test$ export LUA_PATH="../lua/?.lua"
.../UTM/test$ lua -i < MMo.lua
Lua 5.4.7  Copyright (C) 1994-2024 Lua.org, PUC-Rio
> -- test/MMo.lua (RMCG20260409)
>
> mac = require("machine")
> MM = require("MM")
> test = require("test")
>
> -- The 101 machine
> TM101 = mac:new"A0B1> B0C0> C0H1<"
> TM101po = MM.TM2MMo(TM101)
> TM101:compute(true)
101    |A0|->|1B0|->|10C0|->|1H01|
\endcode
As said, when "TM101" starts on a blank tape,
it computes string "101".
We use function
"MM.TM2MMo(tm)" to encode Turing machine "tm",
 including its tape,
into a string that makes a valid initial tape for MMo-UTM.
So "TM101po" is the string that codes
"TM101" with a blank tape for MMo-UTM.
\begincode
> MM.UTMo:tape(TM101po)
MY000X0000111X0101001X1001110Y
> MM.UTMo:steps(900)
1YA1SXM000111X0101001X1001110Y000000000[q45]0   900
\endcode
\new/In \refsc{steps}, we saw
that function\wen/ "tm:steps(n)" runs Turing machine "tm"
until it computes "n" steps or until it halts,
and returns the situation and the number of steps executed.
Notice that MMo-UTM will not halt,
because in state "q45" and reading a "0" ({\em blank}),
it will move to the right without changing the state.
This happens because the emulation has stepped out
the {\em tape} area. However, it works
if we reserve two {\it blank} squares to the right of the scanned square.
\begincode
> TM101xo = string.gsub(TM101po, "M", "M00")
> MM.UTMo:tape(TM101xo)
M00Y000X0000111X0101001X1001110Y
> MM.UTMo:compute()
1M1Y11AXAAAABBBXABABAABXBA01110Y   1103
\endcode
Since the rightmost symbol in the {\em index} is an "A",
then the "M" should be interpreted to be a "0", and
then the emulated {\em tape} is read "101".

\vfil\break

File "test/MMo.lua" also runs
the example in \cite Minsky (1967), §7.3, pages 143--144.
The emulated Turing machine is called "MM73".
Below we show the situations selected by Minsky.
For each situation we show the current state and the tape,
where the scanned square is marked with a bar on top.

The first seven lines reproduce the seven lines
of the unnamed table in page 144,
where the first line is the initial situation,
and the seventh line is the situation after
emulating one step.
The seventh line and the following ones
reproduce Table 7.3-1, in page 143,
which shows the result
after emulating some more steps.

\medskip
\begingroup
\hsize= 10cm
\input ../test/MM73.out
\endgroup

%\vfil\break

\labeled{MMo-Limitations}\subsection{Limitations}

\subsubsection{Unidirectionality}

MMo-UTM only emulates Turing machines with singly left-infinite tapes,
so it will fail to emulate computations that require
more tape to the right than it is available,
as seen in \refsc{MMo-Testing} when emulating \M{TM101}.

Regarding unidirectionality,
\cite Minsky (1967) writes, §6.3.2 in page 129:
\begincitation
It can also be shown that Turing machines which have tapes
infinite in both directions have no advantage or disadvantage
over machines with singly infinite tapes.
\endcitation
He leaves the way of implementing universal Turing machines
without this limitation as an exercise to the reader;
it is problem 7.4-1. In the solution to that problem,
which is in page 286, he makes two proposals:
\beginpoints
\point To {\em fold} the emulated tape, that is,
 to use the squares of the emulated tape alternatively so,
 for example, to emulate the squares to the left we use the odd ones
 and to emulate the squares to the right we use the even ones.
\point To “replace the symbol "M" by a copy of
 the whole instruction region”. Then all the tape would be used
 but, instead of moving one symbol "M",
 all symbols between the two symbols "Y", including them,
 would have to be moved.
\endpoints

\break

His arguments are convincing,
since they help us to see how it could be done,
though they are complex to implement.
Our solution will be easier:
\refx{MMr-UTM}{MMr-UTM} will fulfill the
requirement of bidirectionality by
shifting all the {\em tape} when needed,
see \refsc{MMr-Shift}.
That is, we will implement the algorithm
that makes room for one additional guest in Hilbert's Hotel.

In formal terms, the bidirectionality theorem,
which states that a singly infinite tape
is enough to emulate any Turing machine,
settles this question,
see \refsc{bidirectionality theorem}.


\subsubsection{Binarity}

The other important limitation of MMo-UTM is
that it can only emulate binary Turing machines,
which are those that use a two-symbol alphabet.
\cite Minsky (1967), §6.3 in page 129,
using Shannon's binary theorem, see \refsc{Binary theorem},
writes that
\begincitation
{\em we can restrict our machines to the use of two symbols},
without loss of generality.
\endcitation
In any case, since to eliminate this limitation
requires more work, MMr-UTM will only
emulate binary Turing machines, too.
This limitation will be eliminated later on,
when designing MMc-UTM, in \refsc{MMc-UTM}.


\subsubsection{Dirtiness}

MMo-UTM output string is not clean, it has to be interpreted.
For example,
 when emulating \M{TM101} with enough {\em blank}s to the right,
the output is:
$$\M{1M1Y11AXAAAABBBXABABAABXBA01110Y},$$
which is read \M{101}.
Then the output should rather be \M{101},
and we will design MMr-UTM with the requirement
of cleanliness; see \refsc{MMr-Clean}.


\subsubsection{Starting}

When MMo-UTM starts, the scanned square is the square
where the leftmost symbol \M{X} is written.
Since this is a peculiar requirement,
in the Lua interpreter we had to implement
a \M{resettape} function specific for MMo-UTM.
In order to avoid this anomaly, MMr-UTM will be designed to start from
the leftmost ordinary symbol,
or from any other to its left, which is the usual convention,
see \refsc{situation}.
Although this is not an important requirement,
it is very easy to fulfill, just one additional state,
see \refsc{MMr-Start}.


\subsubsection{Empty table}

Another minor limitation is that
MMo-UTM fails to emulate the empty table Turing machine.
The table is empty when all transitions are halting ones,
so it corresponds to a Turing machine $\cal I$ that does nothing:
 $\forall {\frak d}\in \Gamma^*: \TM I<d> = \TM<d>$.
$\cal I$ is also known as no-operation ("NOP"), and
it implements the literal identity function.
This is an edge case, which was probably overlooked by Minsky,
that is very easy to overcome,
since it does not require additional states,
see \refsc{MMr-Clean}.


\vfil\break

\section{MMr-UTM}

\labeled{MMr-Tape}\subsection{Tape}

A peculiarity of MMo-UTM is that,
after having read the symbol written
in the scanned square of the {\em tape},
it overwrites it with an "M", but taking care of
saving its value in the rightmost square of the {\em index}.
For this reason, when it starts,
 the value of the "M" square should be interpreted
 to be the value of the rightmost bit in the {\em index}.
And it happens the same when it halts.
Therefore, after halting,
 we have to interpret the value of "M" to be "0" or "1"
 depending on the value of the rightmost symbol of the {\em index}
 to be "A" or "B".
To simplify this, we have added another symbol to MMr-UTM,
so the scanned square will be either
 "N", for {\em one}, meaning that emulated scanned square is "1",
 or "M", meaning that emulated scanned square is {\em blank}.
With the additional symbol "N", the alphabet of MMr-UTM is
$\Gamma_{\!\rm r} = \{ \M0, \M1, \M A, \M B, \M M, \M N,
 \M S, \M X, \M Y \}$, where \M0 is the {\em blank} symbol.

An additional requirement is that,
whenever it halts, MMr-UTM will return the very same tape that
the emulated binary Turing machine ${\cal P}_2$ computes.
Therefore, before halting, MMr-UTM has to blank
the {\em program} and {\em index}, and it has to tidy up
the {\em tape}, which requires limiting it.
To this end, MMr-UTM uses another symbol "Y", the third one,
to mark the leftmost position of the {\em tape} that has been reached.
This leftmost "Y" on the {\em tape} is a {\em blank} for ${\cal P}_2$,
and it has to be the leftmost ordinary symbol on the tape of MMr-UTM,
except temporarily while it is being moved to the left.

A valid initial tape for MMr-UTM is as follows,
 where the regular expression uses
 the meta-symbols seen for MMo-UTM, see \refsc{MMo-Tape}.
$$\catcode`\<=1 \catcode`\>=2 \catcode`\{=12 \catcode`\}=12 \tt
\hbox<0*Y[01]*[MN][01]*Y[01]{s}0(X([01]{s}[01][01]{s}[01][01])?)+Y0*>$$

\MT:
\MTbeginfig(15cm,3cm,0pt);
\MTresetpa;
\MT: pickup medpen;
%\MT: draw (0,0) -- (w,0) -- (w,h) -- (0,h) -- cycle;
\MT: numeric cn; cn = 64; numeric ch; ch = 20pt;
\MT: draw (0,h) .. (w,h);
\MT: draw (0,h-ch) .. (w,h-ch);
\MT:
\MT: for n := 1 step 2 until cn:
\MT:  draw (n*w/cn,h) .. (n*w/cn,h-ch); endfor
\MT: string tape;
\MT: tape = "0YM11000Y010X0101110X0110110Y00";
\MT: string sy;
\MT: defaultfont:="cmtt12";
\MT: for n := 0 step 1 until (length tape):
\MT:  sy := substring (n,n+1) of tape;
\MT:  z[n].lbl = ( 2*(n+1)*w/cn, h-ch/2 );
\MT:  label(sy,z[n].lbl);
\MT: endfor
\MT: pickup thinpen;
\MT: defaultfont:="cmti12";
\MT: labeloffset:=3pt;

\MT: draw (17w/cn,h-ch) .. (17w/cn,0);
\MT: drawarrow  (17w/cn,8pt) .. (0,8pt);
\MT: label.top("tape",(8.5w/cn,8pt));

\MT: draw (19w/cn,h-ch) .. (19w/cn,0);
\MT: draw (25w/cn,h-ch) .. (25w/cn,0);
\MT: drawdblarrow  (19w/cn,8pt) .. (25w/cn,8pt);
\MT: label.top("index",(22w/cn,8pt));

\MT: draw (23w/cn,h-ch) .. (23w/cn,(h-ch)/2);
\MT: drawdblarrow  (19w/cn,(h-ch)/2+8pt) .. (23w/cn,(h-ch)/2+8pt);
\MT: label.top("state",(21w/cn,(h-ch)/2+8pt));

\MT: draw (27w/cn,h-ch) .. (27w/cn,(h-ch)/2);
\MT: draw (41w/cn,h-ch) .. (41w/cn,(h-ch)/2);
\MT: drawdblarrow  (27w/cn,(h-ch)/2+8pt) .. (41w/cn,(h-ch)/2+8pt);
\MT: label.top("quintuple",(34w/cn,(h-ch)/2+8pt));

\MT: draw (43w/cn,h-ch) .. (43w/cn,(h-ch)/2);
\MT: draw (57w/cn,h-ch) .. (57w/cn,0pt);
\MT: drawdblarrow  (43w/cn,(h-ch)/2+8pt) .. (57w/cn,(h-ch)/2+8pt);
\MT: label.top("quintuple",(50w/cn,(h-ch)/2+8pt));

\MT: drawdblarrow  (25w/cn,8pt) .. (57w/cn,8pt);
\MT: label.top("program",(41w/cn,8pt));
\MTendfig;

\centerline{\box\MTbox}
%\bigskip

\sssecno=0

\labeled{MMr-Table}\subsection{Table}

\labeled{MMr-Start}\subsubsection{Start}

\MT:
\MTbeginfig(8cm,2cm,0pt);
\MTresetpa;
\MT: pickup medpen;
\MT: defaultfont:="cmtt12";
%\MT: draw (0,0)--(w,0)--(w,h)--(0,h)--cycle;
\MT:  z01 = (1w/4, 1h/2);
\MT:  z11 = (3w/4, 1h/2);
\MTstate(01)"q00"">";
\MTstate(11)"q11""<";
\MT: pickup thinpen;
\MTinput(01.iii)(11.ix)"X""X";
\MTendfig;

Another peculiarity of MMo-UTM is that
it does not start from the leftmost ordinary symbol,
or from a {\em blank} at its left, a convention we follow,
but it starts from the leftmost "X" on the tape,
which is the "X" separating the {\em index} from the {\em program}.
MMr-UTM resolves this problem by adding a new {\em initial} state "q00":
in state "q00", it moves to the right until an "X" is found,
and then it goes to state "q11".
These are the first lines of file "MMr-UTM.csv".
\par\kern-12pt\vbox to0pt{\kern6pt\hbox to 0pt{\kern5cm\box\MTbox\hss}\vss}
\pfile{../csv/MMr-UTM.csv}
-- start
-- locate


\labeled{MMr-Shift}\subsubsection{Shift}

Our \mod[main<last] requirement for MMr-UTM is to avoid the
singly infinite tape limitation of MMo-UTM.
Minsky makes two proposals, see \refsc{Unidirectionality},
but we will implement a simpler solution in MMr-UTM:
to shift all the {\em tape} when needed.
That is, when the emulated movement is going to intrude into
the {\em index}, MMr-UTM will shift every written square
on the {\em tape} one position to the left.
This task also requires the use of the third symbol "Y" that
marks the leftmost position of the tape that has been reached.

In case the emulation tries to move onto the second "Y",
that is, if being MMr-UTM in state "q44", it reads a "Y",
then it keeps moving to the left until finding the leftmost "Y"
but remembering each bit to write it in the next to the left square.
Firstly, state "q51" writes the "M" and,
depending on the symbol overwritten,
it goes to state "q52" or to state "q53".
State "q52" remembers to have read a "0",
and "q53" to have read a "1".
Finally, state "q54" writes the new "Y" and,
since it has read a {\em blank}, it goes to state "q45".
In case the emulation tries to move onto the first "Y",
that is, if being in state "q43", MMr-UTM reads a symbol "Y",
then it overwrites it with an "M", moves to the left
and goes to state "q54", which rewrites the "Y".

\MT:
\MTbeginfig(14cm,8cm,0pt);
\MTresetpa;
\MT: pickup medpen;
\MT: defaultfont:="cmtt12";
%\MT: draw (0,0)--(w,0)--(w,h)--(0,h)--cycle;
\MTcode
 z43 = (1w/6, 1h/4);
 z44 = (5w/6, 1h/4);
 z51 = (4w/6, 2h/4);
 z52 = (3w/6, 3h/4);
 z53 = (3w/6, 1h/4);
 z54 = (2w/6, 2h/4);
 z45 = (1w/6, 3h/4);

\MTstate(43)"q43""<";
\MTstate(44)"q44"">";
\MTstate(51)"q51""<";
\MTstate(52)"q52""<";
\MTstate(53)"q53""<";
\MTstate(54)"q54""<";
\MTstate(45)"q45"">";

\MT: pickup thinpen;

\MTinput(43.ii)(54.vii)"Y""M";
\MTinput(44.x)(51.v)"Y""Y";

\MTinput(51.xi)(52.iv)"0""M";
\MTinput(51.vii)(53.ii)"1""M";
\MTinput(52.vii)(53.xi)"1""0";
\MTinput(52.viii)(54.i)"Y""0";
%\MTauto(52)(i,xi)"0""0";
\MTinput(53.i)(52.v)"0""1";
\MTinput(53.x)(54.v)"Y""1";
%\MTauto(53)(v,vii)"1""1";
\MTinput(54.xi)(45.iv)"0""Y";

\MTendfig;

\par\kern-12pt\vbox to0pt{\kern-14.4pt\hbox to 0pt{\kern3cm\box\MTbox\hss}\vss}
\listline{../csv/MMr-UTM.csv}{73}
\listline{../csv/MMr-UTM.csv}{77}
\pfile{../csv/MMr-UTM.csv}
-- shift
-- clean


\labeled{MMr-Clean}\subsubsection{Clean}

Since MMo-UTM halts when the whole {\em program}
has been inspected without finding any instruction,
then cleaning starts when, in state "q16",
the rightmost symbol "Y" is read.
States "q13" and "q14" only read a "Y" when
 the {\em program} is just an "X",
that is, when all ${\cal P}_2$ transtions are halting ones,
see \refsc{Empty table}.
The cleaning task is simple:
 after state "q16", or "q13" or "q14", has erased the rightmost "Y",
 state \mod["q71"<"q01"] moves to the left blanking the {\em program}
and the {\em index}, and then state \mod["q72"<"q02"] tidies up the {\em tape}
and halts after erasing the leftmost~"Y".
Having added the symbol "N" simplifies this task,
since now the read symbol has not been overwritten.

\MT:
\MTbeginfig(13cm,4cm,0pt);
\MTresetpa;
\MT: pickup medpen;
\MT: defaultfont:="cmtt12";
%\MT: draw (0,0)--(w,0)--(w,h)--(0,h)--cycle;
\MTcode
 z13 = (1w/8, 0.5h/4);
 z14 = (1w/8, 3.5h/4);
 z16 = (1w/8, 2h/4);
 z71 = (3w/8, 2h/4);
 z72 = (5w/8, 2h/4);
 z99 = (7w/8, 2h/4);

\MTstate(13)"q13"">";
\MTstate(14)"q14"">";
\MTstate(16)"q16"">";
\MTstate(71)"q71""<";
\MTstate(72)"q72""<";
\MTstate(99)"q99"">";

\MT: pickup thinpen;

\MTinput(13.iii)(71.viii)"Y""0";
\MTinput(14.iii)(71.x)"Y""0";
\MTinput(16.iii)(71.ix)"Y""0";
\MTinput(71.iii)(72.ix)"Y""0";
\MTauto(71)(vii,v)"*""0";
\MTauto(72)(vii,v)"M""0";
\MTauto(72)(xi,i)"N""1";
\MTinput(72.iii)(99.ix)"Y""0";

\MTendfig;

\par\kern-12pt\vbox to0pt{\kern12pt\hbox to 0pt{\kern3.5cm\box\MTbox\hss}\vss}
\listline{../csv/MMr-UTM.csv}{16}
\new/\wen/\listline{../csv/MMr-UTM.csv}{20}
\listline{../csv/MMr-UTM.csv}{25}
\pfile{../csv/MMr-UTM.csv}
-- clean
-- end


\vfil\break

\labeled{MMr-Patch}\subsubsection{Patch}

And that is all.
This is the patch that
 transforms "MMo-UTM.csv" into "MMr-UTM.csv".
It was produced this way: "diff MMo-UTM.csv MMr-UTM.csv > MMo2MMr.patch".
\doublecolumns
\new/\wen/\file{../csv/MMo2MMr.patch}
\singlecolumn

In the patch we can see the three sections that implement
\refx{Start}{MMr-Start}, \refx{Shift}{MMr-Shift},
 and \refx{Clean}{MMr-Clean},
links to the new sections
 (going to states \mod["q71"<"q01"], "q54", and "q51"),
and then those that deal with symbol "N".
Compared with MMo-UTM, MMr-UTM uses one additional symbol ("N")
and seven states more ("q00", "q51"--"q54", and \mod["q71"<"q01"]--\mod["q72"<"q02"]).

MMr-UTM is a 9-symbol 31-state Turing machine.
One of the 31 states is the halting state \M{q99}.
In the next page, you can find the complete diagram of MMr-UTM.


\labeled{MMr-Testing}\subsection{Testing}

You can use the Turing machine interpreter to test MMr-UTM.
It could be helpful to run
\begincode
.../UTM/test$ export LUA_PATH="../lua/?.lua"
.../UTM/test$ lua -i < MMr.lua > MMr.out
\endcode
and then examining the resulting file ".../UTM/test/MMr.out".

Similarly to MMo-UTM, function "MM.TM2MMr(tm)" encodes
Turing machine "tm", including its tape,
into a string that makes a valid initial tape for MMr-UTM.
And function "MM.MMr2TM(tp)" decodes string "tp",
which is a valid initial MMr-UTM tape,
returning the corresponding Turing machine.

% \medskip
% \file{../test/MMr.out}

\vfill\break

\MT:
\MTbeginfig(15.5cm,23.5cm,0pt); %% MMr-UTM Diagram
\MTresetpa;
\MT: pickup medpen;
\MT: defaultfont:="cmtt12";
\MT: draw (0,0)--(w,0)--(w,h)--(0,h)--cycle withcolor green;
\MTcode
 z09 = (4w/8, 10h/11);
 z01 = (4w/8, 8h/11);
 z02 = (4w/8, 7h/11);
 z99 = (4w/8, 6h/11);
 z11 = (6w/8, 10h/11);
 z12 = (2w/8, 9h/11);
 z13 = (1w/8, 10h/11);
 z14 = (1w/8, 8h/11);
 z15 = (1w/8, 9h/11);
 z16 = (4w/8, 9h/11);
 z21 = (2w/8, 6h/11);
 z22 = (3w/8, 5h/11);
 z23 = (1w/8, 5h/11);
 z24 = (3w/8, 4h/11);
 z25 = (1w/8, 4h/11);
 z26 = (2w/8, 5h/11);
 z31 = (3w/8, 2h/11);
 z32 = (1w/8, 2h/11);
 z33 = (2w/8, 1h/11);
 z34 = (4w/8, 1h/11);
 z35 = (6w/8, 1h/11);
 z41 = (5w/8, 3h/11);
 z42 = (7w/8, 3h/11);
 z43 = (5w/8, 4h/11);
 z44 = (7w/8, 4h/11);
 z45 = (6w/8, 9h/11);
 z46 = (7w/8, 9h/11);
 z51 = (6w/8, 5h/11);
 z52 = (5.4w/8, 6h/11);
 z53 = (6.6w/8, 6h/11);
 z54 = (6w/8, 7h/11);

\MTline{ label("START",(z09+(0,24pt))) withcolor blue;}
\MTstate(09)"q00"">";
\MTstate(01)"q71""<";
\MTstate(02)"q72""<";
\MTstate(99)"q99"">";
\MTline{ label("HALT",(z99-(0,24pt))) withcolor red;}

\MTstate(11)"q11""<";
\MTstate(12)"q12"">";
\MTstate(13)"q13"">";
\MTstate(14)"q14"">";
\MTstate(15)"q15""<";
\MTstate(16)"q16"">";

\MTstate(21)"q21"">";
\MTstate(22)"q22""<";
\MTstate(23)"q23""<";
\MTstate(24)"q24"">";
\MTstate(25)"q25"">";
\MTstate(26)"q26"">";

\MTstate(31)"q31""<";
\MTstate(32)"q32""<";
\MTstate(33)"q33"">";
\MTstate(34)"q34"">";
\MTstate(35)"q35""<";

\MTstate(41)"q41""<";
\MTstate(42)"q42""<";
\MTstate(43)"q43""<";
\MTstate(44)"q44"">";
\MTstate(45)"q45"">";
\MTstate(46)"q46"">";

\MTstate(51)"q51""<";
\MTstate(52)"q52""<";
\MTstate(53)"q53""<";
\MTstate(54)"q54""<";

\MT: pickup thinpen;

\MTinput(09.iii)(11.ix)"X""X";

\MTinput(01.vi)(02.xii)"Y""0";
\MTauto(01)(i,iii)"*""0";
\MTauto(02)(xi,ix)"M""0";
\MTauto(02)(vii,ix)"N""1";
\MTinput(02.vi)(99.xii)"Y""0";

\MTauto(11)(ii,xii)"0""A";
\MTauto(11)(x,xii)"1""B";
\MTinput(11.viii)(12.ii)"Y""Y";
\MTinput(12.xi)(13.v)"A""0";
\MTinput(12.vii)(14.i)"B""1";
\MTinput(12.vi)(21.xii)"X""X";
\MTinput(13.vi)(15.xii)"0""A";
\MTinput(13.iii)(16.x)"1""B";
\MTline{ z13.Y.p = (2.7w/8,9.1h/11);}
\MTinput(13.iv)(01.x)"Y""0";
\MTinput(14.xii)(15.vi)"1""B";
\MTinput(14.iii)(16.viii)"0""A";
\MTinput(14.iv)(01.viii)"Y""0";
\MTinput(15.iii)(12.ix)"Y""Y";
\MTinput(16.ii)(11.vii)"X""X";
\MTinput(16.vi)(01.xii)"Y""0";

\MTinput(21.v)(22.xi)"0""A";
\MTinput(21.vii)(23.i)"1""B";
\MTinput(22.vi)(24.xii)"Y""Y";
\MTinput(23.vi)(25.xii)"Y""Y";
\MTinput(24.xi)(26.iv)"0""A";
\MTinput(24.x)(26.v)"1""A";
\MTinput(24.vi)(31.xii)"X""X";
\MTinput(25.i)(26.viii)"0""B";
\MTinput(25.ii)(26.vii)"1""B";
\MTinput(25.vi)(32.xii)"X""X";
\MTinput(26.xii)(21.vi)"X""X";

\MTinput(31.vii)(33.i)"M""A";
\MTinput(31.v)(33.ii)"N""A";
\MTinput(32.v)(33.xi)"M""B";
\MTinput(32.vii)(33.x)"N""B";
\MTauto(33)(ix,vii)"B""1";
\MTauto(33)(v,vii)"A""0";
\MTinput(33.iii)(34.ix)"X""X";
\MTinput(34.ii)(35.x)"0""0";
\MTinput(34.iv)(35.viii)"1""1";
\MTauto(35)(iii,v)"B""1";
\MTauto(35)(vii,v)"A""0";
\MTinput(35.xi)(41.v)"0""S";
\MTinput(35.i)(42.vii)"1""S";

\MTinput(41.xii)(43.vi)"A""0";
\MTinput(41.ii)(44.viii)"B""0";
\MTinput(42.x)(43.iv)"A""1";
\MTinput(42.xii)(44.vi)"B""1";
\MTline{ z43.O.p = (4.7w/8,6h/11);}
\MTinput(43.xi)(45.vii)"0""M";
\MTline{ z43.I.p = (4.8w/8,6h/11);}
\MTinput(43.xii)(46.vii)"1""N";
\MTline{ z43.Y.p = (4.9w/8,5h/11);}
\MTline{ z43.Y.q = (4.9w/8,6h/11);}
\MTinput(43.i)(54.ix)"Y""M";
\MTline{ z44.O.p = (7.1w/8,6h/11);}
\MTinput(44.xii)(45.v)"0""M";
\MTline{ z44.I.p = (7.3w/8,6h/11);}
\MTinput(44.i)(46.v)"1""N";
\MTinput(44.xi)(51.v)"Y""Y";
\MTinput(45.xii)(11.vi)"S""A";
\MTinput(46.xi)(11.v)"S""B";

\MTinput(51.xi)(52.v)"0""M";
\MTinput(51.i)(53.vii)"1""M";
\MTinput(52.iv)(53.viii)"1""0";
\MTinput(52.i)(54.vii)"Y""0";
\MTinput(53.x)(52.ii)"0""1";
\MTinput(53.xi)(54.v)"Y""1";
\MTinput(54.xii)(45.vi)"0""Y";

\MTline{ label("MMr-UTM",(4w/8,12pt));}

\MTendfig;

\dest
\new/\wen/\centerline{\define*|MMrDiagram|MMr-UTM|MMo-UTM
 revised, see \refsc{MMr-UTM}|\box\MTbox}

\vfil\break

\section{Binary Turing completeness theorem}

Whenever it halts,
MMr-UTM returns the very same binary string that
the emulated binary Turing machine ${\cal P}_2$ computes. And,
if ${\cal P}_2$ does not halt, then MMr-UTM will not halt, either.
Therefore, MMr-UTM is a perfect emulator of binary Turing machines.
The {\em theorem of binary Turing completeness}\define*|binary Turing completeness theorem||%
there is a Turing machine that can emulate any two-symbol Turing machine| says:\\
{\sc Theorem}: {\em There is a binary-universal Turing machine}.

{\sc Proof}: In other words,
$$\exists\,{\cal U}_2\in{\frak T},
  \forall {\cal P}_2\in{\frak T}_2,
  \forall {\frak d}_2\in\Gamma_2^* \;:\;
  {\cal U}_2\langle\,{\frak p}_2| {\frak d}_2\,\rangle
   =  {\cal P}_2\langle{\frak d}_2\rangle$$
where:
\beginpoints
\point $\frak T$ is the set of Turing machines,
 so ${\cal U}_2$ is a Turing machine,
\point ${\frak T}_2$ is the set of binary Turing machines,
 which are those that use only two symbols,
 {\em blank}, coded \M0, and another symbol, coded \M1,
 so their alphabet is
 $\Gamma_{\!2} = \{ \M0, \M1 \}$,
 and then ${\cal P}_2$ is a binary Turing machine
 and ${\frak d}_2$ is a binary string,
\point ${\cal P}_2\langle{\frak d}_2\rangle \in \Gamma^*_2$
 is the output string computed by ${\cal P}_2$ when
 its input is ${\frak d}_2$,\\
 or $\infty\notin\Gamma^*_2$ if it does not halt,
\point ${\frak p}_2\in {\cal L}_{{\cal U}_2}$
 is the program that encodes ${\cal P}_2$\\
 in the language ${\cal L}_{{\cal U}_2}$ implemented by ${\cal U}_2$,
 ${\cal L}_{{\cal U}_2} \subset \Gamma^*_{\!{\cal U}_2}$,
\point the $|$ represents the composing operation of
 language ${\cal L}_{{\cal U}_2}$, so
 ${\frak p}_2| {\frak d}_2 \in  {\cal L}_{{\cal U}_2}$, and
\point ${\cal U}_2\langle\,{\frak p}_2| {\frak d}_2\,\rangle
 \in \Gamma^*_2$ is the corresponding output
 string computed by ${\cal U}_2$,\\
 or $\infty\notin\Gamma^*_2$ if it does not halt.
\endpoints

Since MMr-UTM, in the place of ${\cal U}_2$,
satisfies the equation,
we have shown that the theorem is {\sc true}.
\QED

MMr-UTM is a constructive proof of the theorem.
Notice that MMo-UTM does not satisfy the equation
because of its unidirectionality, see \refsc{Unidirectionality},
because of its dirtiness, see \refsc{Dirtiness},
and because it fails to emulate $\cal I$, see \refsc{Empty table}.

\label{Binary theorem}
The only remaining limitation of MMr-UTM is that the emulated
Turing machine has to be binary, that is, it has to use
only two symbols, {\em blank} and another one.
Regarding this limitation \cite Shannon (1956), page 163, wrote:
\begincitation
It is also possible, as we will now show,
to construct a machine, C,
which will act like any given Turing machine A
and use only two symbols 1 and 0 on its tape,
one of which, 0 say, is the symbol for a blank square.
\endcitation
Here, we will call it the
\define|binary theorem||every Turing machine
 can be emulated by a two-symbol one|
by Shannon.

Following Shannon, to emulate an arbitrary
$m$-symbol $n$-state Turing machine A
with a universal Turing machine that only emulates
two-symbol Turing machines, as MMr-UTM,
we would have to compute the corresponding
binary Turing machine C, which, by his calculations,
has at most $3n2^w+n(2w-7)$ states,
where $w\in{\bb N}$ is such that $2^{w-1} < m \le 2^w$.
This way, we would be transferring complexity
from the universal Turing machine to the encoding,
see \refsc{Discussion}.

Together, the binary Turing completeness theorem and
Shannon's binary theorem
show that there is a universal Turing machine.
So next we will design MMc-UTM, a universal Turing machine,
that is, one that can emulate any Turing machine,
without exceptions. We will base MMc-UTM on  MMr-UTM,
which is MMo-UTM revised.


\vfil\break

\section{MMc-UTM}

\labeled{MMc-Tape}\subsection{Tape}

In order to emulate any Turing machine $\cal P$,
MMc-UTM uses the same encoding for programs and data as MMo-UTM,
with minor differences, as follows.
MMo-UTM requires a binary fixed-length encoding
for the emulated states $\Phi_{\cal P}$, and MMc-UTM too.
We are denoting the number of bits that encode an emulated state "s".
MMo-UTM uses a binary fixed-length {\em index}, and MMc-UTM too,
which thus requires a binary fixed-lengh encoding
for the emulated alphabet $\Gamma_{\!\cal P}$.
We will denote the number of bits that encode an emulated symbol "w".
Another requirement is that
 the emulated {\em blank} is a string of {\em blank}s,
that is, a string of "w" symbols "0" of MMc-UTM.
Aditionally, in order to avoid ambiguities,
the leftmost bit of all other emulated symbols should be "1",
though this is not a requirement.
Finally, MMo-UTM codes emulated {\em left} \M0
and emulated {\em right} \M1, and MMc-UTM too.

The initial tape of Misky's MMo-UTM has exactly one "M" on the tape,
marking the emulated scanned square.
The initial tape of MMc-UTM has "w" consecutive "M" or "N" symbols,
where "w" in the number of bits needed to encode each emulated symbol.
Save for this, the tape of MMc-UTM is organized the same way that
the tape of MMr-UTM, see \refsc{MMr-Tape}.
As MMr-UTM, compared with MMo-UTM,
MMc-UTM uses an additional symbol "N",
so the alphabet of MMc-UTM is
 $\Gamma_{\!\rm c} = \{ \M0, \M1, \M A, \M B, \M M, \M N,
 \M S, \M X, \M Y \}$, where "0" is the {\em blank} symbol.

The initial tape of MMc-UTM contains
 symbol "Y" three times,
 "w" times in total symbols "M" and "N",
 symbol "X" as many times as quintuples in $\cal P$, and at least one,
 symbol "1" a finite number of times,
 and any number of {\em blank}s, "0".
The initial tape of MMc-UTM is any described by the regular expression
\begincode
0*Y[01]*[MN]{w}[01]*Y[01]{s}[01]{w}(X([01]{s}[01]{w}[01]{s}[01]{w}[01])?)+Y0*
\endcode
which uses the meta-symbols seen for MMo-UTM, see \refsc{MMo-Tape}.

The tape of MMc-UTM has three areas:
\beginpoints
\point {\em tape}: emulated tape of $\cal P$,
 where the consecutive symbols "M" and "N" mark
 the emulated scanned square and represent its value,
 "M" for "0" and "N" for "1",
 and where the leftmost~"Y", which represents a "0",
 indicates that all the symbols to its left are "0";
\point {\em index}: current/next state ({\em st}) and
 read/write symbol ({\em sy}) of $\cal P$;
\point {\em program}: description of $\cal P$,
 that is, its linearized table encoded in fixed length binary,
 done by prefixing each binary encoded quintuple of $\cal P$
 with a symbol "X".
\endpoints

\MT:
\MTbeginfig(15cm,3cm,0pt);
\MT: pickup medpen;
%\MT: draw (0,0) -- (w,0) -- (w,h) -- (0,h) -- cycle;
\MT: numeric cn; cn = 64; numeric ch; ch = 20pt;
\MT: draw (0,h) .. (w,h);
\MT: draw (0,h-ch) .. (w,h-ch);
\MT:
\MT: for n := 1 step 2 until cn:
\MT:  draw (n*w/cn,h) .. (n*w/cn,h-ch); endfor
\MT: string tape;
\MT: tape = "YMN10Y0001X000011100X000111010Y";
\MT: string sy;
\MT: defaultfont:="cmtt12";
\MT: for n := 0 step 1 until (length tape):
\MT:  sy := substring (n,n+1) of tape;
\MT:  z[n].lbl = ( 2*(n+1)*w/cn, h-ch/2 );
\MT:  label(sy,z[n].lbl);
\MT: endfor
\MT: pickup thinpen;
\MT: defaultfont:="cmti12";
\MT: labeloffset:=3pt;

\MT: draw (11w/cn,h-ch) .. (11w/cn,0);
\MT: drawarrow  (11w/cn,8pt) .. (0,8pt);
\MT: label.top("tape",(6w/cn,8pt));

\MT: draw (13w/cn,h-ch) .. (13w/cn,0);
\MT: draw (21w/cn,h-ch) .. (21w/cn,0);
\MT: drawdblarrow  (13w/cn,8pt) .. (21w/cn,8pt);
\MT: label.top("index",(17w/cn,8pt));

\MT: draw (17w/cn,h-ch) .. (17w/cn,(h-ch)/2);
\MT: drawdblarrow  (13w/cn,(h-ch)/2+8pt) .. (17w/cn,(h-ch)/2+8pt);
\MT: label.top("st",(15w/cn,(h-ch)/2+8pt));
\MT: drawdblarrow  (17w/cn,(h-ch)/2+8pt) .. (21w/cn,(h-ch)/2+8pt);
\MT: label.top("sy",(19w/cn,(h-ch)/2+8pt));

\MT: draw (23w/cn,h-ch) .. (23w/cn,(h-ch)/2);
\MT: draw (41w/cn,h-ch) .. (41w/cn,(h-ch)/2);
\MT: drawdblarrow  (23w/cn,(h-ch)/2+8pt) .. (41w/cn,(h-ch)/2+8pt);
\MT: label.top("quintuple",(32w/cn,(h-ch)/2+8pt));

\MT: draw (43w/cn,h-ch) .. (43w/cn,(h-ch)/2);
\MT: draw (61w/cn,h-ch) .. (61w/cn,0pt);
\MT: drawdblarrow  (43w/cn,(h-ch)/2+8pt) .. (61w/cn,(h-ch)/2+8pt);
\MT: label.top("quintuple",(52w/cn,(h-ch)/2+8pt));

\MT: drawdblarrow  (21w/cn,8pt) .. (61w/cn,8pt);
\MT: label.top("program",(41w/cn,8pt));

\MTendfig;

\bigskip
\centerline{\box\MTbox}
\medskip

We will follow the emulation of Turing machine "TM1022",
 see \refsc{TM1022b}, by MMc-UTM.
The initial tape for our "TM1022" example is shown below,
where the scanned square is marked with a bar on top, which is,
initially, the leftmost non-{\em blank} square:
\smallskip\continuepar
\hbox{\qquad\tt\=YMMY0000X000001101X001000100X010010111X011011000X011100110X100001110Y}


\vfill\break

\openin\situations=../test/pr1022.out


\labeled{MMc-Table}\subsection{Table}

\sssecno=-1

\labeled{MMc-Start}\subsubsection{Start}


% MMc-UTM has several sub-machines,
% which we will explain one by one.

\MT:
\MTbeginfig(8cm,2cm,0pt);
\MTresetpa;
\MT: pickup medpen;
\MT: defaultfont:="cmtt12";
%\MT: draw (0,0)--(w,0)--(w,h)--(0,h)--cycle;
\MT:  z01 = (1w/4, 1h/2);
\MT:  z11 = (3w/4, 1h/2);
\MTstate(01)"q00"">";
\MTstate(11)"q11""<";
\MT: pickup thinpen;
\MTinput(01.iii)(11.ix)"X""X";
\MTendfig;

The initial situation of our "TM1022" example is:

\nextsituation
%\tape{q00}{\=YMMY0000X000001101X001000100X010010111X011011000X011100110X100001110Y}

This first sub-machine,
which we take from MMr-UTM, see \refsc{MMr-Start},
only uses one state, the {\em initial} state "q00",
which moves to the right until the leftmost "X" is found.
Once the proper square is pointed to,
it goes to state "q11" to locate the instruction,
see \refsc{MMc-Locate}.
So these are the first lines of file "MMc-UTM.csv".
\par\kern-12pt\vbox to0pt{\kern14.4pt\hbox to 0pt{\kern5cm\box\MTbox\hss}\vss}
\pfile{../csv/MMc-UTM.csv}
-- start
-- locate

When this sub-machine terminates the first time, the situation is as follows.

\nextsituation
%\tape{q11}{YMMY000\=0X000001101X001000100X010010111X011011000X011100110X100001110Y}


\labeled{MMc-Locate}\subsubsection{Locate}

We take the locate sub-machine of MMo-UTM.
As \cite Minsky (1967), page 139, explains:
\begincitation \font\tt=cmsltt10 at 12pt
It first searches to the right to find the first
state-symbol pair that matches the one represented in
the [{\em index}] area. On the path to the desired
pair, it changes all "0"'s and "1"'s to "A"'s and "B"'s.
After finding the desired pair,
and changing it also to "A"'s and "B"'s,
it runs back to the leftmost "X".
\endcitation

\MT:
\MTbeginfig(12cm,8cm,0pt);
\MTresetpa;
\MT: pickup medpen;
\MT: defaultfont:="cmtt12";
%\MT: draw (0,0)--(w,0)--(w,h)--(0,h)--cycle;

\MTcode
 z70 = (1w/8, 7h/9);
 z11 = (3w/8, 7h/9);
 z12 = (5w/8, 4h/9);
 z13 = (6w/8, 7h/9);
 z14 = (6w/8, 1h/9);
 z15 = (7w/8, 4h/9);
 z16 = (3w/8, 4h/9);
 z21 = (3w/8, 1h/9);
 z71 = (1w/8, 4h/9);

\MTstate(70)"q00"">";
\MTstate(11)"q11""<";
\MTstate(12)"q12"">";
\MTstate(13)"q13"">";
\MTstate(14)"q14"">";
\MTstate(15)"q15""<";
\MTstate(16)"q16"">";
\MTstate(21)"q21"">";
\MTstate(71)"q71""<";

\MT: pickup thinpen;

\MTinput(70.iii)(11.ix)"X""X";
\MTauto(11)(xi,i)"0""A";
\MTauto(11)(iii,i)"1""B";
\MTinput(11.v)(12.x)"Y""Y";
\MTinput(12.i)(13.vii)"A""0";
\MTinput(12.v)(14.xi)"B""1";
\MTinput(12.vii)(21.ii)"X""X";
\MTinput(13.v)(15.xi)"0""A";
\MTinput(13.viii)(16.ii)"1""B";
\MTinput(13.x)(71.ii)"Y""0";
\MTinput(14.i)(15.vii)"1""B";
\MTinput(14.x)(16.iv)"0""A";
%\MTline{ z14.Y.p = (3.5w/8, 2h/9); }
\MTinput(14.viii)(71.iv)"Y""0";
\MTinput(15.ix)(12.iii)"Y""Y";
\MTinput(16.xii)(11.vi)"X""X";
\MTinput(16.ix)(71.iii)"Y""0";

\MTendfig;

If an instruction is found,
then it goes to state "q21" to copy the instruction
 from the {\em program} into de {\em index}, see \refsc{MMc-Copy}.
If no instruction is found,
then it goes to state \mod["q71"<"q01"] to finish, see \refsc{MMc-Clean};
 as done also for MMr-UTM, see \refsc{MMr-Clean},
 we have added two transitions to state \mod["q71"<"q01"],
one from "q13" and the other form "q14",
 in order to deal with empty tables.
\par\kern-12pt\vbox to0pt{\kern24pt\hbox to 0pt{\kern4cm\box\MTbox\hss}\vss}
\new/\wen/\pfile{../csv/MMc-UTM.csv}
-- locate
-- copy

When this sub-machine terminates the first time,
 the situation is as follows.

\nextsituation
%\tape{q21}{YMMY0000X\=AAAA01101X001000100X010010111X011011000X011100110X100001110Y}


\labeled{MMc-Copy}\subsubsection{Copy}

We take without any modification the copy sub-machine of MMo-UTM.
As \cite Minsky (1967), pages 139 and 140, explains:
\begincitation \font\tt=cmsltt10 at 12pt
It starts on the leftmost "X", and from there shifts to the right
 until it has gone past the last "A" or "B" and
 sees for the first time some some "0"'s or "1"'s.
 These "0"'s or "1"'s describe the new state of [$\cal P$],
 the new symbol which must eventually be printed in location "M"["N"],
 and the direction of motion.
 Next the machine copies the "0"'s and "1"'s representing
 the [new state] and [new symbol] into the [{\em index}] region;
 it does not print [the direction of motion] but just remembers it.
\endcitation


\MT:
\MTbeginfig(12cm,7.5cm,0pt);
\MTresetpa;
\MT: pickup medpen;
\MT: defaultfont:="cmtt12";
%\MT: draw (0,0)--(w,0)--(w,h)--(0,h)--cycle;
\MTcode
 z12 = (1w/8, 5h/6);
 z21 = (4w/8, 5h/6);
 z22 = (2w/8, 3h/6);
 z23 = (6w/8, 3h/6);
 z24 = (3w/8, 1h/6);
 z25 = (5w/8, 1h/6);
 z26 = (4w/8, 3h/6);
 z31 = (1w/8, 1h/6);
 z35 = (7w/8, 1h/6);

\MTstate(12)"q12"">";
\MTstate(21)"q21"">";
\MTstate(22)"q22""<";
\MTstate(23)"q23""<";
\MTstate(24)"q24"">";
\MTstate(25)"q25"">";
\MTstate(26)"q26"">";
\MTstate(31)"q31""<";
\MTstate(35)"q35""<";

\MT: pickup thinpen;

\MTinput(12.iii)(21.ix)"X""X";
\MTinput(21.vii)(22.i)"0""A";
\MTinput(21.v)(23.xi)"1""B";
\MTinput(22.v)(24.xi)"Y""Y";
\MTinput(23.vii)(25.i)"Y""Y";
\MTinput(24.i)(26.viii)"0""A";
\MTinput(24.ii)(26.vii)"1""A";
\MTinput(24.ix)(31.iii)"X""X";
\MTinput(25.x)(26.v)"1""B";
\MTinput(25.xi)(26.iv)"0""B";
\MTinput(25.iii)(35.ix)"X""X";
\MTinput(26.xii)(21.vi)"X""X";

\MTendfig;


It copies the instruction from the {\em program} to the {\em index},
both the next state and the symbol to write.
State "q21" moves to the right till
 the leftmost unread bit in the {\em program},
and if its a "0" then it changes it to "A" and goes to state "q22",
but if its a "1" then it changes it to "B" and goes to state "q23".
State "q22" (and "q23") moves to the left until the middle "Y",
and then changes to state "q24" (to "q25").
State "q24" (and "q25") moves to the right till
 the leftmost unwritten bit in the {\em index},
it changes it to "A" (to "B"), and goes to state "q26".
State "q26" moves to the leftmost "X", which marks the beginning
of the {\em program}, and goes to state "q21" to copy the next bit.
If in state "q24" (or "q25") there are not any unwritten bit
 in the {\em index}, then it is the movement,
so it goes to state "q31" to emulate the left movement
(or to "q35" to emulate the right movement), see \refsc{MMc-Move}.
\par\kern-12pt\vbox to0pt{\kern14.4pt\hbox to 0pt{\kern4cm\box\MTbox\hss}\vss}
\pfile{../csv/MMc-UTM.csv}
-- copy
-- move left

When this sub-machine terminates the first time,
since the emulated movement is {\em right},
the situation is as follows.

\nextsituation
%\tape{q35}{YMMYABB\=AXAAAAABBABX001000100X010010111X011011000X011100110X100001110Y}


\vfill\break

\labeled{MMc-Move}\subsubsection{Move}

It changes as many "0"'s or "1"'s to "A"'s and "B"'s
as there are "M"'s and "N"'s. This way later, when reading,
these marked bits will be found among all the "0"'s and "1"'s
on the {\em tape}.
In the process, the "M"'s and "N"'s are changed to "S"'s,
so they will be the target when writing.


\MT:
\MTbeginfig(12cm,6cm,0pt);
\MTresetpa;
\MT: pickup medpen;
\MT: defaultfont:="cmtt12";
%\MT: draw (0,0)--(w,0)--(w,h)--(0,h)--cycle;
\MTcode
 z24 = (1w/8, 3h/4);
 z31 = (3w/8, 3h/4);
 z32 = (5w/8, 3h/4);
 z33 = (5w/8, 1h/4);
 z34 = (7w/8, 2h/4);
 z51 = (2w/8, 1h/4);
 z52 = (2w/8, 2h/4);

\MTstate(24)"q24"">";
\MTstate(31)"q31""<";
\MTstate(32)"q32""<";
\MTstate(33)"q33"">";
\MTstate(34)"q34""<";
\MTstate(51)"q51"">";
\MTstate(52)"q52""<";

\MT: pickup thinpen;

\MTinput(24.iii)(31.ix)"X""X";
\MTinput(31.ii)(32.x)"M""S";
\MTinput(31.iv)(32.viii)"N""S";
\MTinput(32.vii)(33.xi)"0""A";
\MTinput(32.v)(33.i)"1""B";
\MTinput(32.iii)(34.x)"Y""A";
\MTinput(33.ix)(51.iii)"0""0";
\MTinput(33.viii)(51.iv)"1""1";
\MTinput(33.xii)(32.vi)"M""S";
\MTline{ z33.N.p = (5.7w/8,2h/4);}
\MTinput(33.ii)(32.iv)"N""S";
%\MTinput(33.v)(32.vii)"N""S";
\MTinput(33.x)(52.iv)"Y""Y";
\MTinput(34.viii)(33.iii)"0""Y";

\MTendfig;


Let us see first the movement to the left.
In this case, the "A"'s and "B"'s are written to the left
of the "M"'s and "N"'s by state "q32".
But firstly, state "q31" has moved from the {\em program},
 where a left movement was read, through the {\em index},
 until the rightmost "M" or "N" on the {\em tape},
 and then another square to the left.
If state "q32" finds a "Y",
then the inspected {\em tape} has to be enlarged to the left,
which is done by state "q34".
State "q33" determines whether to go back to "q32" to continue marking,
or go forward to  state "q51", or to "q52", to write, see \refsc{MMc-Write}.
\par\kern-12pt\vbox to0pt{\kern12pt\hbox to 0pt{\kern4cm\box\MTbox\hss}\vss}
\pfile{../csv/MMc-UTM.csv}
-- move left
-- move right

\MT:
\MTbeginfig(12cm,6cm,0pt);
\MTresetpa;
\MT: pickup medpen;
\MT: defaultfont:="cmtt12";
%\MT: draw (0,0)--(w,0)--(w,h)--(0,h)--cycle;
\MTcode
 z25 = (1w/8, 3h/4);
 z35 = (3w/8, 3h/4);
 z36 = (5w/8, 3h/4);
 z37 = (5w/8, 1h/4);
 z41 = (7w/8, 3h/4);
 z48 = (7w/8, 1h/4);
 z51 = (2w/8, 1h/4);

\MTstate(25)"q25"">";
\MTstate(35)"q35""<";
\MTstate(36)"q36"">";
\MTstate(37)"q37""<";
\MTstate(41)"q41""<";
\MTstate(48)"q48"">";
\MTstate(51)"q51"">";

\MT: pickup thinpen;

\MTinput(25.iii)(35.ix)"X""X";
\MTinput(35.ii)(36.x)"M""S";
\MTinput(35.iv)(36.viii)"N""S";
\MTinput(36.vii)(37.xi)"0""A";
\MTinput(36.v)(37.i)"1""B";
\MTinput(36.iii)(41.ix)"Y""Y";
\MTinput(37.x)(51.ii)"0""0";
\MTinput(37.viii)(51.iv)"1""1";
\MTinput(37.xii)(36.vi)"M""S";
\MTline{ z37.N.p = (5.7w/8,2h/4);}
\MTinput(37.ii)(36.iv)"N""S";
\MTinput(37.ix)(51.iii)"Y""Y";
\MTinput(48.ix)(37.iii)"S""S";

\MTendfig;

Let us now see the movement to the right.
In this case, the "A"'s and "B"'s are written to the right
of the "M"'s and "N"'s by state "q36".
But firstly, state "q35" has moved from the {\em program},
 where a right movement was read,
through the {\em index}, until the rightmost "M" or "N" on the {\em tape}.
Then state "q36" marks the bits to be read,
but, if it finds the "Y" separating the {\em tape} from the {\em index},
then it goes to state "q41" to make room, see \refsc{MMc-Shift}.
Finally, state "q37" determines whether to go back to "q36"
to continue marking,
or go forward to state "q51" to write, see \refsc{MMc-Write}.
\par\kern-12pt\vbox to0pt{\kern6pt\hbox to 0pt{\kern4cm\box\MTbox\hss}\vss}
\pfile{../csv/MMc-UTM.csv}
-- move right
-- shift

When this sub-machine terminates the first time,
since the emulated movement was {\em right} and
it finds a "Y", the situation is as follows.

\nextsituation
%\tape{q41}{YM\=TYABBAXAAAAABBABX001000100X010010111X011011000X011100110X100001110Y}


\vfill\break

\labeled{MMc-Shift}\subsubsection{Shift}

\MT:
\MTbeginfig(12cm,15cm,0pt);
\MTresetpa;
\MT: pickup medpen;
\MT: defaultfont:="cmtt12";
%\MT: draw (0,0)--(w,0)--(w,h)--(0,h)--cycle;
\MTcode
 z36 = (1w/8, 7h/8);
 z37 = (1w/8, 1h/8);
 z41 = (4w/8, 7h/8);
 z42 = (7w/8, 7h/8);
 z43 = (4w/8, 5h/8);
 z44 = (7w/8, 5h/8);
 z45 = (4w/8, 3h/8);
 z46 = (7w/8, 3h/8);
 z47 = (7w/8, 1h/8);
 z48 = (4w/8, 1h/8);

\MTstate(36)"q36"">";
\MTstate(37)"q37""<";
\MTstate(41)"q41""<";
\MTstate(42)"q42""<";
\MTstate(43)"q43""<";
\MTstate(44)"q44""<";
\MTstate(45)"q45""<";
\MTstate(46)"q46""<";
\MTstate(47)"q47""<";
\MTstate(48)"q48"">";

\MT: pickup thinpen;

\MTinput(36.iii)(41.ix)"Y""Y";
\MTinput(41.ii)(42.x)"B""A";
\MTinput(41.vi)(43.xii)"M""A";
\MTinput(41.vii)(43.xi)"N""A";
\MTinput(41.v)(44.xi)"S""A";
\MTinput(42.ix)(41.iii)"A""B";
\MTinput(42.vii)(43.ii)"M""B";
\MTinput(42.viii)(43.i)"N""B";
\MTinput(42.vi)(44.xii)"S""B";
\MTinput(43.vi)(45.xii)"0""M";
\MTinput(43.v)(46.xi)"1""M";
\MTauto(43)(viii,x)"N""M";
\MTinput(43.iii)(44.ix)"S""M";
\MTline{ z43.Y.p = (3w/8,3h/8);}
\MTline{ z43.Y.q = (4w/8,2h/8);}
\MTline{ z43.Y.t = (5w/8,1.67h/8);}
\MTinput(43.vii)(47.x)"Y""M";
\MTinput(44.vii)(45.i)"0""S";
\MTinput(44.vi)(46.xii)"1""S";
\MTinput(44.viii)(43.iv)"M""S";
\MTinput(44.x)(43.ii)"N""S";
\MTline{ z44.Y.p = (8w/8,3h/8);}
\MTinput(44.v)(47.ii)"Y""S";

\MTinput(45.ii)(46.x)"1""0";
\MTinput(45.v)(47.xi)"Y""0";
\MTinput(46.viii)(45.iv)"0""1";
\MTinput(46.vi)(47.xii)"Y""1";
\MTinput(47.ix)(48.iii)"0""Y";
\MTinput(48.ix)(37.iii)"S""S";

\MTendfig;

If trying to move to the right, it finds the middle "Y",
that is, if the emulation is trying to intrude into the {\em index},
then it has to make room for the symbol,
so it will shift all the symbols on the tape.
Conceptually, it is part of the right movement;
it comes from state "q36" when there is no room,
and goes back to the state "q37" to complete the movement.
Though it has many transitions, is simple.
States "q41"--"q46" do the main shifting,
each one remembering the symbol read, "A", "B", "M", "S", "0" or "1".
State "q47" writes the leftmost "Y" one square to the left,
and state "q48" moves to the leftmost "S" and
goes back to state "q37" to complete the movement,
 see \refsc{MMc-Move}.
\par\kern-12pt\vbox to0pt{\kern-6pt\hbox to 0pt{\kern4cm\box\MTbox\hss}\vss}
\pfile{../csv/MMc-UTM.csv}
-- shift
-- write

When this sub-machine terminates the first time,
the situation is as follows.

\nextsituation
%\tape{q37}{Y\=MTAYABBAXAAAAABBABX001000100X010010111X011011000X011100110X100001110Y}

And when the first emulated movement is completed,
the situation is as follows.

\nextsituation
%\tape{q51}{Y\=TTAAYABBAXAAAAABBABX001000100X010010111X011011000X011100110X100001110Y}


\vfill\break

\labeled{MMc-Write}\subsubsection{Write}

\MT:
\MTbeginfig(12cm,9cm,0pt);
\MTresetpa;
\MT: pickup medpen;
\MT: defaultfont:="cmtt12";
%\MT: draw (0,0)--(w,0)--(w,h)--(0,h)--cycle;
\MTcode
 z33 = (1w/8, 7h/8);
 z37 = (5w/8, 7h/8);
 z51 = (3w/8, 7h/8);
 z52 = (1w/8, 5h/8);
 z53 = (3w/8, 5h/8);
 z54 = (1w/8, 2h/8);
 z55 = (3w/8, 3h/8);
 z56 = (3w/8, 1h/8);
 z57 = (5w/8, 2h/8);
 z61 = (7w/8, 2h/8);

\MTstate(33)"q33"">";
\MTstate(37)"q37""<";
\MTstate(51)"q51"">";
\MTstate(52)"q52""<";
\MTstate(53)"q53"">";
\MTstate(54)"q54""<";
\MTstate(55)"q55""<";
\MTstate(56)"q56""<";
\MTstate(57)"q57""<";
\MTstate(61)"q61"">";

\MT: pickup thinpen;

\MTinput(33.ii)(51.x)"0""0";
\MTinput(33.iv)(51.viii)"1""1";
\MTinput(33.vi)(52.xii)"Y""Y";
\MTinput(37.x)(51.ii)"0""0";
\MTinput(37.viii)(51.iv)"1""1";
\MTinput(37.ix)(51.iii)"Y""Y";

\MTinput(51.vii)(52.i)"Y""Y";
\MTinput(52.iii)(53.ix)"S""M";
\MTinput(53.vii)(54.i)"X""X";
\MTinput(54.ii)(55.ix)"A""S";
\MTinput(54.iv)(56.ix)"B""S";
\MTinput(55.iii)(57.x)"M""0";
\MTinput(56.iii)(57.viii)"M""1";

\MTinput(57.xi)(53.v)"S""M";
\MTinput(57.i)(61.xi)"0""0";
\MTinput(57.ii)(61.x)"1""1";
\MTinput(57.iii)(61.ix)"A""A";
\MTinput(57.iv)(61.viii)"B""B";
\MTinput(57.v)(61.vii)"Y""Y";

\MTendfig;

Every symbol "S" on the {\em tape} is overwritten with
the symbol that is in the {\em index}.
State "q51" moves to the "Y" that separates the {\em tape}
from the {\em index}, and state "q52" moves to the rightmost "S",
and changes it to an "M". Then
state "q53" moves to the "X" that separates the {\em index}
from the {\em program}, and state "q54" moves to the rightmost
"A" or "B", which overwrites with an "S",
and it remembers it going to state "q55" or "q56".
States "q55" and "q56" move to the "M" and overwrite it
accordingly, that is, "q55" with a "0" and "q56" with a "1".
Then, state "q57" inspects the square at its left;
if it is a "S", it overwrites it with an "M" and goes to "q53"
to write another bit; otherwise it goes to state "q61"
to read the symbol, see \refsc{MMc-Read}.
\par\kern-12pt\vbox to0pt{\kern6pt\hbox to 0pt{\kern4cm\box\MTbox\hss}\vss}
\pfile{../csv/MMc-UTM.csv}
-- write
-- read

When this sub-machine terminates the first time,
the situation is as follows.

\nextsituation
%\tape{q61}{Y\=10AAYABSSXAAAAABBABX001000100X010010111X011011000X011100110X100001110Y}


\vfill\break

\labeled{MMc-Read}\subsubsection{Read}

\MT:
\MTbeginfig(9cm,8cm,0pt);
\MTresetpa;
\MT: pickup medpen;
\MT: defaultfont:="cmtt12";
%\MT: draw (0,0)--(w,0)--(w,h)--(0,h)--cycle;
\MTcode
 z57 = (5w/6, 5h/6);
 z61 = (1w/6, 5h/6);
 z62 = (1w/6, 3h/6);
 z63 = (3w/6, 3h/6);
 z64 = (3w/6, 1h/6);
 z65 = (1w/6, 1h/6);
 z66 = (5w/6, 1h/6);
 z69 = (5w/6, 3h/6);

\MTstate(57)"q57""<";
\MTstate(61)"q61"">";
\MTstate(62)"q62"">";
\MTstate(63)"q63""<";
\MTstate(64)"q64""<";
\MTstate(65)"q65"">";
\MTstate(66)"q66"">";
\MTstate(69)"q00"">";

\MT: pickup thinpen;

\MTinput(57.xi)(61.i)"0""0";
\MTinput(57.x)(61.ii)"1""1";
\MTinput(57.ix)(61.iii)"A""A";
\MTinput(57.viii)(61.iv)"B""B";
\MTinput(57.vii)(61.v)"Y""Y";

\MTinput(61.vi)(62.xii)"Y""Y";
\MTinput(62.ii)(63.x)"0""0";
\MTinput(62.iv)(63.viii)"1""1";
\MTauto(62)(xi,ix)"A""0";
\MTauto(62)(vii,ix)"B""1";
\MTinput(62.iii)(63.ix)"Y""Y";
\MTinput(63.vi)(64.xii)"S""X";
\MTinput(63.iii)(69.ix)"Y""Y";
\MTinput(64.ix)(65.iii)"A""M";
\MTinput(64.iii)(66.ix)"B""N";
\MTinput(65.i)(63.vii)"X""0";
\MTinput(66.xi)(63.v)"X""1";

\MTendfig;

The bits marked on the {\em tape} are
writen in the symbol part of the {\em index}.
State "q61" moves to the middle "Y".
State "q62" moves right through the {\em index} and {\em program}
restoring every "A" to "0" and every "B" to "1" until
finding the first "0" or "1", or the rightmost "Y".
State "q63" moves to the rightmost "S" that overwrites with an "X"
and goes to state "q64".
State "q64" moves to the rightmost "A" or "B",
and remembers it going to state "q65" or to "q66".
State "q65" then moves to the "X" and overwrites it with a "0";
and state "q66" with a "1".
Both go back to "q63" to read another bit.
However, if in state "q63" there is not any remaining "S",
the reading task is completed, the emulated step finishes,
and MMc-UTM goes back to the initial state "q00",
see \refsc{MMc-Start}.
\par\kern-12pt\vbox to0pt{\kern12pt\hbox to 0pt{\kern5cm\box\MTbox\hss}\vss}
\pfile{../csv/MMc-UTM.csv}
-- read
-- clean

When this sub-machine terminates the first time,
the first emulated step is completed, and the situation is as follows.

\nextsituation
%\tape{q00}{Y10MMY\=0100X000001101X001000100X010010111X011011000X011100110X100001110Y}

\continuepar
This is a valid initial situation, save for the scanned square,
which is not the leftmost "Y".

\vfill\break

\labeled{MMc-Clean}\subsubsection{Clean}

We take without any modification the clean sub-machine of MMr-UTM,
see \refsc{MMr-Clean}.
When the computation has been completed,
the only reamining task is to clean the tape,
by erasing the {\em program} and the {\em index}
and tidying up the {\em tape}.
This is the situation.

\skipsituation{7}
\new/\wen/\nextsituation
%\tape{q71}{YNM001111Y11BAXAAAAABBABXAABAAABAAXABAABABBBXABBABBAAAXABBBAABBAXBA000111\=00}

\MT:
\MTbeginfig(13cm,4cm,0pt);
\MTresetpa;
\MT: pickup medpen;
\MT: defaultfont:="cmtt12";
%\MT: draw (0,0)--(w,0)--(w,h)--(0,h)--cycle;
\MTcode
 z13 = (1w/8, 0.5h/4);
 z14 = (1w/8, 3.5h/4);
 z16 = (1w/8, 2h/4);
 z01 = (3w/8, 2h/4);
 z02 = (5w/8, 2h/4);
 z99 = (7w/8, 2h/4);

\MTstate(13)"q13"">";
\MTstate(14)"q14"">";
\MTstate(16)"q16"">";
\MTstate(01)"q71""<";
\MTstate(02)"q72""<";
\MTstate(99)"q99"">";

\MT: pickup thinpen;

\MTinput(13.iii)(01.viii)"Y""0";
\MTinput(14.iii)(01.x)"Y""0";
\MTinput(16.iii)(01.ix)"Y""0";
\MTinput(01.iii)(02.ix)"Y""0";
\MTauto(01)(vii,v)"*""0";
\MTauto(02)(vii,v)"M""0";
\MTauto(02)(xi,i)"N""1";
\MTinput(02.iii)(99.ix)"Y""0";

\MTendfig;

Since MMc-UTM detects a {\em halt} instruction when the whole {\em program}
has been inspected without finding any instruction,
then cleaning starts when, in state "q16", or "q13" or "q14",
the rightmost symbol "Y" is read, which is erased.
Next it goes to state \mod["q71"<"q01"], which erases
everything until the middle "Y", that is,
it blanks the {\em program} and the {\em index}.
And then it goes to state \mod["q72"<"q02"], which tidies up the {\em tape},
and it finally halts after erasing the leftmost "Y".
\par\kern0pt\vbox to0pt{\kern-15pt\hbox to 0pt{\kern3.75cm\box\MTbox\hss}\vss}
\new/\wen/\pfile{../csv/MMc-UTM.csv}
-- clean
-- end

\bigskip

This is the situation when the halting state is reached.

\skipsituation{1}
\nextsituation
%\\tape{q99}{0\=10001111000000000000000000000000000000000000000000000000000000000000000000}

\closein\situations

\continuepar
Disregarding the infinities of {\em blank}s,
the output is "10001111" which decodes as "1022".

The advantage of encoding every ordinary symbol
with a leftmost "1", is that then
the leftmost non-{\em blank} bit of the emulation output is
the leftmost bit of the leftmost ordinary symbol.
In other words, this way
the emulation is always as ambiguous as the emulated computation.

On the other hand, if, for example, we would have encoded
 Turing machine "TM1022" applying the map
$$\{ \M{0}\leftrightarrow\M{00},
     \M{1}\leftrightarrow\M{01},
     \M{2}\leftrightarrow\M{10} \}\;,$$
then the output of MMc-UTM emulating it on a blank tape
would have been
$$\dots\M0\M0\M0\M1\M0\M0\M1\M0\M1\M0\M0\M0\dots\,,$$
which is ambiguous,
 either "01-00-10-10" decoded as "1022"
 or "10-01-01" decoded as "211",
while the emulated computation is not ambiguous, is "1022".

We say that {\em blank} is a {\em special} symbol
because of these infinities of {\em blank}s,
which are always present on the tapes of Turing machines,
and which are prone to cause ambiguities when concatenated
with finite strings, or substrings, of {\em blank} symbols.

\bigskip
\cleartabs
MMc-UTM is a 9-symbol 44-state Turing machine.
In the next page, you can find the complete diagram of MMc-UTM.
\+ Sub-machines &
 \refx{Locate}{MMc-Locate} (\M{q1n}) and
 \refx{Copy}{MMc-Copy} (\M{q2n}),
 with \del/\refx{Start}{MMc-Start} and\led/\refx{Clean}{MMc-Clean} (\M{q7n}),
 & use \mod[14<15] states,\cr
\+& \refx{Move}{MMc-Move} (\M{q3n}) and
 \refx{Shift}{MMc-Shift} (\M{q4n}) & use 15 states,\cr
\+& \refx{Write}{MMc-Write} (\M{q5n}) and
 \refx{Read}{MMc-Read} (\M{q6n}) & use 13 states,\cr
\continuepar
which, adding \new/the initial state \M{q00} and\wen/ the halting state \M{q99}, make 44 states.\cr

\vfill\break

\MT:
\MTbeginfig(15.5cm,23.5cm,0pt); %% MMc-UTM Diagram
\MTresetpa;
\MT: pickup medpen;
\MT: defaultfont:="cmtt12";
\MT: draw (0,0)--(w,0)--(w,h)--(0,h)--cycle withcolor green;
\MTcode
 z09 = (4w/8, 10h/11);
 z71 = (4w/8, 8h/11);
 z72 = (4w/8, 7h/11);
 z99 = (4w/8, 6h/11);
 z11 = (3w/8, 10h/11);
 z12 = (2w/8, 9h/11);
 z13 = (1w/8, 10h/11);
 z14 = (1w/8, 8h/11);
 z15 = (1w/8, 9h/11);
 z16 = (4w/8, 9h/11);
 z21 = (2w/8, 7h/11);
 z22 = (3w/8, 6h/11);
 z23 = (1w/8, 6h/11);
 z24 = (3w/8, 5h/11);
 z25 = (1w/8, 5h/11);
 z26 = (2w/8, 6h/11);
 z31 = (3w/8, 4h/11);
 z32 = (4w/8, 4h/11);
 z33 = (6w/8, 4h/11);
 z34 = (5w/8, 5h/11);
 z35 = (1w/8, 4h/11);
 z36 = (1w/8, 3h/11);
 z37 = (7w/8, 3h/11);
 z41 = (1w/8, 2h/11);
 z42 = (1w/8, 1h/11);
 z43 = (3w/8, 2h/11);
 z44 = (3w/8, 1h/11);
 z45 = (5w/8, 2h/11);
 z46 = (5w/8, 1h/11);
 z47 = (7w/8, 1h/11);
 z48 = (7w/8, 2h/11);
 z51 = (7w/8, 4h/11);
 z52 = (7w/8, 5h/11);
 z53 = (6w/8, 5.5h/11);
 z54 = (5w/8, 7h/11);
 z55 = (6w/8, 7.5h/11);
 z56 = (6w/8, 6.5h/11);
 z57 = (7w/8, 7h/11);
 z61 = (7w/8, 8h/11);
 z62 = (7w/8, 9h/11);
 z63 = (6w/8, 9h/11);
 z64 = (6w/8, 10h/11);
 z65 = (5w/8, 10h/11);
 z66 = (7w/8, 10h/11);

\MTline{ label("START",(z09+(0,24pt))) withcolor blue;}
\MTstate(09)"q00"">";
\MTstate(71)"q71""<";
\MTstate(72)"q72""<";
\MTstate(99)"q99"">";
\MTline{ label("HALT",(z99-(0,24pt))) withcolor red;}

\MTstate(11)"q11""<";
\MTstate(12)"q12"">";
\MTstate(13)"q13"">";
\MTstate(14)"q14"">";
\MTstate(15)"q15""<";
\MTstate(16)"q16"">";

\MTstate(21)"q21"">";
\MTstate(22)"q22""<";
\MTstate(23)"q23""<";
\MTstate(24)"q24"">";
\MTstate(25)"q25"">";
\MTstate(26)"q26"">";

\MTstate(31)"q31""<";
\MTstate(32)"q32""<";
\MTstate(33)"q33"">";
\MTstate(34)"q34""<";
\MTstate(35)"q35""<";
\MTstate(36)"q36"">";
\MTstate(37)"q37""<";

\MTstate(41)"q41""<";
\MTstate(42)"q42""<";
\MTstate(43)"q43""<";
\MTstate(44)"q44""<";
\MTstate(45)"q45""<";
\MTstate(46)"q46""<";
\MTstate(47)"q47""<";
\MTstate(48)"q48"">";

\MTstate(51)"q51"">";
\MTstate(52)"q52""<";
\MTstate(53)"q53"">";
\MTstate(54)"q54""<";
\MTstate(55)"q55""<";
\MTstate(56)"q56""<";
\MTstate(57)"q57""<";

\MTstate(61)"q61"">";
\MTstate(62)"q62"">";
\MTstate(63)"q63""<";
\MTstate(64)"q64""<";
\MTstate(65)"q65"">";
\MTstate(66)"q66"">";

\MT: pickup thinpen;

\MTinput(09.ix)(11.iii)"X""X";

\MTauto(11)(ii,xii)"0""A";
\MTauto(11)(x,xii)"1""B";
\MTinput(11.vii)(12.i)"Y""Y";
\MTinput(12.xi)(13.v)"A""0";
\MTinput(12.vii)(14.i)"B""1";
\MTinput(12.vi)(21.xii)"X""X";
\MTinput(13.vi)(15.xii)"0""A";
\MTinput(13.iii)(16.x)"1""B";
\MTline{ z13.Y.p = (2.7w/8,9.1h/11);}
\MTinput(13.iv)(71.x)"Y""0";
\MTinput(14.xii)(15.vi)"1""B";
\MTinput(14.iii)(16.viii)"0""A";
\MTinput(14.iv)(71.viii)"Y""0";
\MTinput(15.iii)(12.ix)"Y""Y";
\MTinput(16.xi)(11.v)"X""X";
\MTinput(16.vi)(71.xii)"Y""0";

\MTinput(21.v)(22.xi)"0""A";
\MTinput(21.vii)(23.i)"1""B";
\MTinput(22.vi)(24.xii)"Y""Y";
\MTinput(23.vi)(25.xii)"Y""Y";
\MTinput(24.x)(26.v)"0""A";
\MTinput(24.xi)(26.iv)"1""A";
\MTinput(24.vi)(31.xii)"X""X";
\MTinput(25.i)(26.viii)"1""B";
\MTinput(25.ii)(26.vii)"0""B";
\MTinput(25.vi)(35.xii)"X""X";
\MTinput(26.xii)(21.vi)"X""X";

\MTinput(31.ii)(32.x)"M""S";
\MTinput(31.iv)(32.viii)"N""S";
\MTinput(32.ii)(33.x)"0""A";
\MTinput(32.iii)(33.ix)"1""B";
\MTinput(32.i)(34.vii)"Y""A";
\MTinput(33.ii)(51.x)"0""0";
\MTinput(33.iv)(51.viii)"1""1";
\MTinput(33.viii)(32.iv)"M""S";
\MTinput(33.vii)(32.v)"N""S";
\MTinput(33.i)(52.vii)"Y""Y";
\MTinput(34.v)(33.xi)"0""Y";
\MTinput(35.vii)(36.xi)"M""S";
\MTinput(35.v)(36.i)"N""S";

\MTinput(36.ii)(37.x)"0""A";
\MTinput(36.iii)(37.ix)"1""B";
\MTinput(36.vi)(41.xii)"Y""Y";
\MTinput(37.xi)(51.vii)"0""0";
\MTinput(37.xii)(51.vi)"1""1";
\MTinput(37.viii)(36.iv)"M""S";
\MTinput(37.vii)(36.v)"N""S";
\MTinput(37.i)(51.v)"Y""Y";

\MTinput(41.vi)(42.xii)"B""A";
\MTinput(41.iii)(43.ix)"M""A";
\MTinput(41.iv)(43.viii)"N""A";
\MTinput(41.v)(44.x)"S""A";
\MTinput(42.xi)(41.vii)"A""B";
\MTinput(42.i)(43.viii)"M""B";
\MTinput(42.ii)(43.vii)"N""B";
\MTinput(42.iii)(44.ix)"S""B";
\MTinput(43.iii)(45.ix)"0""M";
\MTinput(43.iv)(46.x)"1""M";
\MTauto(43)(x,xii)"N""M";
\MTinput(43.vi)(44.xii)"S""M";
\MTline{ z43.Y.p = (5w/8,2.5h/11);}
\MTinput(43.ii)(47.xi)"Y""M";
\MTinput(44.ii)(45.viii)"0""S";
\MTinput(44.iii)(46.ix)"1""S";
\MTinput(44.xi)(43.vii)"M""S";
\MTinput(44.i)(43.v)"N""S";
\MTline{ z44.Y.p = (5w/8,0.5h/11);}
\MTinput(44.iv)(47.viii)"Y""S";

\MTinput(45.vii)(46.xi)"1""0";
\MTinput(45.iv)(47.x)"Y""0";
\MTinput(46.i)(45.v)"0""1";
\MTinput(46.iii)(47.ix)"Y""1";
\MTinput(47.xii)(48.vi)"0""Y";
\MTinput(48.xii)(37.vi)"S""S";

\MTinput(51.xii)(52.vi)"Y""Y";
\MTinput(52.x)(53.iv)"S""M";
\MTinput(53.xi)(54.vi)"X""X";
\MTinput(54.ii)(55.viii)"A""S";
\MTinput(54.iv)(56.x)"B""S";
\MTinput(55.iv)(57.x)"M""0";
\MTinput(56.ii)(57.viii)"M""1";

\MTinput(57.vi)(53.i)"S""M";
\MTinput(57.xi)(61.vii)"0""0";
\MTinput(57.xii)(61.vi)"1""1";
\MTinput(57.i)(61.v)"A""A";
\MTline{ z57.B.p = (7.4w/8,7.5h/11);}
\MTinput(57.ii)(61.iv)"B""B";
\MTline{ z57.Y.p = (7.6w/8,7.5h/11);}
\MTinput(57.iii)(61.iii)"Y""Y";

\MTinput(61.xii)(62.vi)"Y""Y";
\MTinput(62.x)(63.ii)"0""0";
\MTinput(62.viii)(63.iv)"1""1";
\MTauto(62)(v,iii)"A""0";
\MTauto(62)(i,iii)"B""1";
\MTinput(62.ix)(63.iii)"Y""Y";
\MTinput(63.xii)(64.vi)"S""X";
\MTinput(63.x)(09.iv)"Y""Y";
\MTinput(64.viii)(65.iv)"A""M";
\MTinput(64.iv)(66.viii)"B""N";
\MTinput(65.v)(63.xi)"X""0";
\MTinput(66.vii)(63.i)"X""1";

\MTinput(71.vi)(72.xii)"Y""0";
\MTauto(71)(ii,iv)"*""0";
\MTauto(72)(xi,ix)"M""0";
\MTauto(72)(vii,ix)"N""1";
\MTinput(72.vi)(99.xii)"Y""0";

\MTline{ label("MMc-UTM",(4w/8,12pt));}

\MTendfig;

\dest
\new/\wen/\centerline{\define*|MMcDiagram|MMc-UTM|universal Turing
 machine based on MMo-UTM, see \refsc{MMc-UTM}|\box\MTbox}

\vfill\break


\labeled{MMc-Testing}\subsection{Testing}

We repeat here the emulation of Turing machine "TM1022" by MMc-UTM.
The first line shows the initial situation.
Thereafter, each time state "q00" is reached,
 MMc-UTM has completed a emulated step.
It takes MMc-UTM 7397 steps to emulate the 7 steps of "TM1022",
or of "TM1022b", on a blank tape.
The last line shows the halting situation.
\medskip
\new/\wen/\input ../test/pr1022.out
\medskip

You can use again the Turing machine interpreter to test MMc-UTM.
\begincode
.../UTM/test$ export LUA_PATH="../lua/?.lua"
.../UTM/test$ lua -i < MMc.lua > MMc.out
\endcode
Test suite "test/MMc.lua" includes two auto-emulations,
that is, MMc-UTM emulating itself emulating a Turing machine\mod[.<,]
\mod[The  shortest one takes<that take] \new/$981\,682\,311$ steps,
that is,\wen/ about a billion\mod[.<($10^9$) steps each]
\new/And MMc-UTM emulating itself emulating "TM1022" on a blank tape
 takes $8\,857\,405\,567$ steps, which are nearly nine billions.\wen/
\mod[Therefore<, so then], depending on your machine, it could take a long time to complete.
Then you can examine the resulting output file ".../UTM/test/MMc.out".

Function "MM.TM2MMc(tm)" encodes Turing machine "tm",
 including its tape,
into a string that makes a valid initial tape for MMc-UTM,
and function "MM.MMc2TM(tp)" decodes string "tp",
an initial MMc-UTM tape,
returning the corresponding Turing machine.

Function "MM.recode(s,dic,rpat,fpat,rsep,fsep)" takes
 a string "s" representing a table and returns
 another string representing the same table recoded.
 Firstly, it uses patterns "rpat" and "fpat" to extract
 its records and fields,
 next it codes each field using dictionary "dic",
 and then it reassembles back the
 coded fields into records, using "fsep" as separator,
 and the resulting records into a string,
 using "rsep" as separator.

\del/Function "tm:steps(n,...)"\ executes steps
 until some condition is met.
If Turing machine "tm" halts, then the function returns.
It also returns after having executed "n" steps.
Each of the other arguments ("...")\ refers either
to a state ("s", a "string"),
and then the function returns if the current state is "s",
or  to a transition ("{st,sy}", a state and symbol "table"),
and then the function returns
 if the current state is "st" and the scanned symbol is "sy".\led/


\vfill\break

\section{Turing completeness theorem}

Whenever it halts,
MMc-UTM returns, simply encoded, the same string that
the emulated Turing machine $\cal P$ computes. And,
if $\cal P$ does not halt, then MMc-UTM will not halt, either.
Therefore, MMc-UTM is a perfect emulator of Turing machines.
The {\em theorem of Turing completeness}\define*|Turing-completeness-theorem|Turing
 completeness theorem|there is a universal Turing machine| says:\\
{\sc Theorem}: {\em There is a universal Turing machine}.

\dest
{\sc Proof}: In other words,
$$\exists\,{\cal U}\in{\frak T},\,
  \forall{\cal P}\in{\frak T},\,
  \forall{\frak d}\in\Gamma^*_{\!\cal P}:\;
 \overleftarrow{\TM U<\,p|\vec{d}\,>} = \TM P<d> \;,$$
where
\beginpoints
\point $\frak T$ is the set of Turing machines,
 so $\cal U$ and $\cal P$ are Turing machines,
\point ${\frak d}$ is an input string to $\cal P$,
\point $\TM P<d> \in\Gamma^*_{\!\cal P}$ is the corresponding output
 string, or $\infty \notin \Gamma^*_{\!\cal P}$ if it does not halt,
\point the right pointing arrow hat represents encoding, so
 $\vec{\frak d} \in {\cal L_U}$ is the string $\frak d$
 encoded in the language $\cal L_U$ implemented by $\cal U$,
 ${\cal L_U} \subset \Gamma^*_{\!\cal U}$,
\point ${\frak p} \in {\cal L_U}$ is the
 \define|program||an encoded Turing machine,
 ${\frak p} = \smash{\vec{\cal P}}$|
 that codes $\cal P$ in the language $\cal L_U$,
 so \smash{$\frak p = \vec{\cal P}$},
\point the $|$ represents the composing operation
 of language $\cal L_U$, so $\frak p|\vec{d} \in {\cal L_U}$,
\point $\TM U<\,p|\vec{d}\,> \in {\cal L_U}$ is the
 corresponding output string computed by $\cal U$,\\
 %which is encoded in the language $\cal L_U$,
 or $\infty \notin \Gamma^*_{\!\cal U}$ if it does not halt, and
\point the left pointing arrow on top denotes the corresponding decoding,
 $\forall {\frak d} \in \Gamma^*_{\!\cal P}:\!\!\smash{\overleftarrow{\frak\vec d}} = {\frak d}$,
 which leaves $\infty$ unchanged, $\overleftarrow{\vphantom{i}\infty} = \infty$.
\endpoints

Since MMc-UTM, in the place of $\cal U$, satisfies the equation,
then we have shown that the theorem is {\sc true}.
\QED


Therefore,
MMc-UTM proves constructively
 the theorem of Turing completeness, and
MMc-UTM is a universal Turing machine~$\cal U$.

\dest
By the way,
since it never moves to the right of the rightmost "Y",
MMc-UTM also proves constructively
the {\em theorem of bidirectionality}\define*|bidirectionality theorem||a
 singly infinite tape is enough to emulate any Turing machine|:
{\em a singly infinite tape is enough to emulate any Turing machine}.
In particular, a singly infinite tape universal Turing machine,
such as MMc-UTM, can emulate Turing machines that work on tapes that
are infinite in both directions.


\dest
The language \mod[implemented by<of] a universal Turing machine
$\cal L_U$ is a \define|Turing complete language||the language
 of a universal Turing machine, $\cal L_U$|\mod[.<because]
\new/In a Turing complete language,\wen/
every Turing machine and every computation
can be meaningfully expressed\del/<in it\led/;
program $\frak p$ means Turing machine $\cal P$,
and sentence $\frak p|\vec d$ means computation $\TM P<d>$.

It seems imposible that a complex Turing machine
can be emulated by a simpler universal Turing machine.
However, this is the case. % For example,
MMc-UTM is a 9-symbol 44-state Turing machine that is universal,
and then, even if $\cal P$ is a $\gamma$-symbol $\phi$-state
Turing machine such that $\gamma > 9$ and $\phi > 44$,
it can be emulated by MMc-UTM.



\vfil\break

\section{Discussion}

\subsection{Turing}

The first universal computing machine was designed
by Alan \cite Turing (1936), §7, who presented a UTM*,
which we will call ATo-UTM. In his words (page 241),
ATo-UTM “can be used to compute any computable sequence”,
because his aim was to compute the binary expansion
of any computable number, which is one that
a Turing machine can calculate with arbitrary precision,
see \cite Casares (N).
Since mathematically the binary expansion is as valid as any other,
the binary limitation was not a limitation for Turing.
In other words, sometimes encoding can be abstracted away.

The requirement of arbitrary precision
means that, whenever a bit of the computable number
that is being calculated is commited,
it will not be recalculated anymore. Then, for Turing,
it was enough to compute {\em sequences}, for which
a singly infinite tape suffices. Thus, basically, ATo-UTM has the
same limitations that MMo-UTM has, as seen in \refsc{MMo-Limitations}:
unidirectionality, binarity, and dirtiness.
For details on ATo-UTM limitations, see \cite Petzold (2008), page 161.

Even though ATo-UTM and MMo-UTM have the same limitations,
ATo-UTM is much more complex than MMo-UTM.
The reason, we can guess, is that ATo-UTM was the very first one,
while MMo-UTM had the advantage of examining all previous trials.
The reader interested on ATo-UTM
will find \cite Davies (2004) explanation
clearer than the original exposition by Turing,
for the same reason.
By the way, ATo-UTM is so complex that Davies,
although his aim was just to correct the errors,
had to redesign ATo-UTM in order to run it.
Following our naming practice,
we should call the version revised by Davies ATr-UTM.


\subsection{Encoding}

Sometimes encoding can be abstracted away,
but some other times encoding could be abused.
This problem was addressed by \cite Davis (1956), page 167.
\begincitation
 The universality of a Turing machine is manifested
 by its ability, {\em via a suitable encoding},
 to perform any computation which could be performed
 by any given Turing machine. However, the condition
 must surely be added that the encoding itself be,
 in some suitable sense, simple.
 For there would not be much point in claiming
 universality for a Turing machine
 for which the encoding would require, in essence,
 another universal machine to carry it out.
 This raises the problem of explicitly
 defining {\em universal Turing machine},
 which problem is the subject of the present note.
\endcitation
Then, following Turing,
a universal Turing machine is a Turing machine
that can perform any computation.
And, in addition, Davis requires (page 169) “that
the encoding [of the computation] may be entirely accomplished
by Turing machines which are not themselves universal.”

\dest\cite Minsky (1967),
whole first paragraph of §14.8.1 in page 277,
goes further.
In order to make “sure that the machine was not applied to the data
during the encoding procces”, he requires that
 “{\em the encodings for the machine structure and for the data
 [to be] done separately}.”
This is why, in the equation of Turing completeness,
 see \refsc{Turing completeness theorem},
 program $\frak p$ is Turing machine $\cal P$ encoded and
 argument $\frak\vec d$ is data $\frak d$ encoded {\em separately},
 which are then composed, denoted by $|$.
We will call Minsky's requirement the
\define|separation requirement||program and data
 have to be encoded separately|.

Minsky made those comments because
he had designed a 4-symbol 7-state Turing machine
that can emulate any 2-tag system.
Then, using a theorem by Cocke, presented in page 270,
that affirms that, for any Turing machine $\cal P$,
 there is a 2-tag system that behaves like $\cal P$,
he concluded that his 4-7-machine must be universal,
though the encoding is not straightforward. As he writes,
in page 278, “the conversion [doing the encoding]
is a primitive-recursive operation
[thus fulfilling Davis' requirement]. In fact,
it involves only writing down sixteen productions
for each quintuple.”
On the other hand, Minsky's didactic encoding used by MMo-UTM
is straightforward, just a finite mapping, see \refsc{MMo-UTM}.
And so they are MMr-UTM and MMc-UTM encodings,
 see \refsc{MMr-UTM} and \refsc{MMc-UTM}.

Minsky is right writing, again in page 277, that
“there is not particular virtue in the quintuple formulation,
and one might be able to get a simpler universal Turing machine
with some other representation.”
This has allowed him, and others following his track,
to design small universal Turing machines but,
by requiring an additional theorem,
these universal Turing machines do not make
self-contained proofs of the theorem of Turing completeness.


\subsection{Definition}

Our definition of universal Turing machine
is not anymore the one by Turing (UTM*).
The current definition of universal Turing machine (UTM)
derives from the separation requirement by Minsky.
\beginpoints
\point Turing: a UTM* is a Turing machine that
can perform any computation.
\point Minsky: a UTM is a Turing machine that
 can emulate any Turing machine.
\endpoints
Neither ATo-UTM nor MMo-UTM qualifies as universal Turing machine
under the current definition,
since neither can emulate itself. On the other hand,
any universal Turing machine under the current definition
can perform any computation, and then it
“can be used to compute any computable sequence”,
and then it is also universal by Turing's definition.
In summary, every UTM is a UTM*, but not the converse.

Minsky's separation requirement is the reason why
the current definition of universal Turing machine
has superseded the original one by Turing. Let us see it.

Any computation can be performed in many ways.
Let us take, for example, the computation that calculates
the hundreth bit of the binary expansion of number $\pi$.
There are several algorithms to calculate it,
and for each algorithm several Turing machines that implement it,
and then, at least, one way of encoding each Turing machine
as a program for each universal Turing machine.
Since number $\pi$ is computable,
our computation example is a halting one,
so it will take a finite number of steps
and it will scan a finite number of tape squares.
Then, a singly infinite tape is enough,
meaning that there are programs that
compute it on a singly infinite tape.
In summary, we can compute it using
a UTM* that only emulates singly infinite tapes,
but, since not every program goes,
we have to make some judgements on the behavior of programs.
For example, since MMo-UTM's {\em tape} is {\em only} infinite
to the left, we have to determine in advance
how many squares to the right does each program require.
And this is undecidable,
given the halting problem and its consequences,
in this case \cite[\1's (\2)] Rice (1953) theorem.

The current definition of universal Turing machine
avoids this problem. And, in the case of MMc-UTM,
 which satisfies the current definition,
encoding simply applies a finite mapping to translate
each element (state, symbol or movement) in the Turing machine table
into a binary string, so it is independent of its behavior;
all is syntax, no semantics are involved.


\subsection{Foundations}

In order to compare this proof of the theorem,
in sections \refsc{MMc-UTM} and \refsc{Turing completeness theorem},
with others, it could be instructive to divert ourselves for a moment
to the theory of recursive functions.
In a very precise sense,
recursion theory and computing theory are the same,
since \cite Turing (1937) proved that
every recursive function is computable, and the converse.
In that sense, the theorem of Turing completeness of computing
is equivalent to the normal form theorem of recursion.
However, \cite Kleene (1952) presents his normal form theorem
in page 288, using Gödel numbers and dense mathematical tools.
Of course, it can be done in less pages;
\cite Davis (1958), who also uses recursion theory and Gödel numbers,
presents the theorem of Turing completeness in page 64
and Kleene normal form theorem in page 171.

In stark contrast to those proofs,
the proof in this paper should be easy to follow,
since the universal Turing machine presented, MMc-UTM,
is completely defined,
so it can be run in any Turing machine interpreter,
and it is composed of eight sub-machines,
each one with a specific task accomplished
with eight states, or less.
Once you understand how each sub-machine works,
then it will also be easy to see how the whole machine works.
At that point, you will convince yourself that
there is a universal Turing machine.
Surely, if you design your own universal Turing machine,
which I encorage you to do,
you will feel more assured, but
even in that case you can elaborate on this one, MMc-UTM,
to design yours, as I did on the universal Turing machine
by \cite Minsky (1967), MMo-UTM.
In the end,
the purpose of any mathematical proof should be
to convince its reader that the theorem is {\sc true}.

The advantage of \cite Turing (1936) computing
 compared with other equivalent forms,
such as recursion or lambda-calculus,
is its conceptual simplicity.
For example, it is documented, see \cite Davis (1982),
that, for stating what is effectively calculable,
Gödel disregarded his own recursion
and Church's lambda-calculus, and that
he was instead convinced that Turing computing
was the definitive answer
only after he had read % \cite[\1's (\2)] Turing (1936) analysis,
Turing's analysis,
which is axiomatic, see \cite Shagrir (2006).
Using an analogy, Turing computing is machine code, while
all other equivalent formulations are higher level languages.

Then, an advantage of the proof in this paper
is that, being of the lowest level,
it does not need any abstraction.
Another advantage is that,
neither depending on other theorems nor
on equivalences with other formulations,
it is self-contained.
And together, concreteness and independency,
show the foundational character
of the proof of the theorem of Turing completeness
given in this paper, sections
\refsc{MMc-UTM} and \refsc{Turing completeness theorem}.


\section{Conclusion}

In this paper, we have presented
the whole table of a universal Turing machine
that has not any limitation,
that is conceptually simple, and
that uses a straightforward encoding.
These are qualities that make it
a good foundational and constructive proof of
the theorem of Turing completeness.


\vfill\break

\definitions

\vfill\break

\references
\item Casares (N)
\item Copeland (2004)
\xitem Davies (2004) = Davies (2004x)
\xitem Davis (1956) = Davis (1956x)
\item Davis (1958)
\item Davis (1982)
\item Kleene (1952)
\item Minsky (1967)
\item Petzold (2008)
\item Rice (1953)
\item Shagrir (2006)
\xitem Shannon (1956) = Shannon (1956x)
\item Shannon \& McCarthy (1956)
\item Turing (1936)
\item Turing (1937)
\endreferences


\bye
